% !TEX root = ../高等概率论.tex
% =====================================================================
%  随机变量：可测映射、乘积 σ-代数与截口、不可数族的上确界、生成的 σ-代数、简单函数逼近、函数形式的单调类定理、Doob–Dynkin 引理
% =====================================================================
\lecture{可测映射与随机变量，可测性的判别与运算，乘积 \texorpdfstring{$\sigma$}{σ}-代数与截口，广义实值随机变量，不可数个随机变量的上确界与经验分布函数，由随机变量生成的 \texorpdfstring{$\sigma$}{σ}-代数，简单函数逼近，函数形式的单调类定理，Doob--Dynkin 引理}
\label{lec:random-variables}

初等概率论把随机变量说成是“取值随机的量”。这个名字其实有些误导：随机变量既不随机，也不是变量，它是从样本空间 $\Omega$ 到 $\R$ 的一个完全确定的函数 $X$。随机性全部来自 $\omega$ 的选取：先按照概率 $P$ 随机地选出一个样本点 $\omega$，然后读出数值 $X(\omega)$。所以同一个样本空间上可以同时定义许多随机变量，比如说把硬币抛 $n$ 次，第一次的结果、正面的总数、最长的连续正面的长度，都是同一个 $\omega$ 的函数。

对这个函数，我们只提一个要求：对每个 Borel 集 $B$，“$X$ 落在 $B$ 中”必须是一个事件，否则我们就无法谈论 $P(X\in B)$。这个要求只涉及 $\sigma$-代数 $\cF$，与概率测度 $P$ 无关。所以这一讲的全部内容都不涉及 $P$，它研究的是可测函数这一类对象本身：怎样判断一个函数可测，可测函数经过哪些运算仍然可测。结论是，分析中能够遇到的一切运算，包括四则运算、复合连续函数、取上下确界与极限，都不会破坏可测性。

\topic{本讲用到的前文结果}

\begin{enumerate}
\item 生成的 $\sigma$-代数与好集原理（定义~\ref{def:set-systems:generated} 及其后的讨论）：要证 $\sigma(\cC)$ 中每个集合有某性质，只需证有该性质的集合构成 $\sigma$-代数并包含 $\cC$。
\item $\cR$ 由 $\{(-\infty,x]:x\in\Q\}$ 生成（命题~\ref{prop:set-systems:borel-generators}）；$\cR^d$ 由 $\{(-\infty,x_1]\times\cdots\times(-\infty,x_d]\}$ 生成（习题~\exref{ex:set-systems:A5}）。
\item 原像与生成的 $\sigma$-代数可交换：$\sigma(f^{-1}(\cC))=f^{-1}(\sigma(\cC))$（习题~\exref{ex:set-systems:A7}）。
\item $\pi$-$\lambda$ 定理（定理~\ref{thm:set-systems:pi-lambda}）。
\item 集合列与数列的上、下极限（命题~\ref{prop:set-systems:limsup}）。
\item 第~\ref{lec:extension-of-measures} 讲的 Lebesgue 可测集与不可测集（作业 D），只在注记和习题中用到。
\end{enumerate}

\topic{可测映射与随机变量}

对映射 $X:\Omega\to S$ 与 $B\subseteq S$，记 $X^{-1}(B)=\{\omega\in\Omega: X(\omega)\in B\}$，简写为 $\{X\in B\}$。原像与集合运算可交换：
\[
X^{-1}(B^c)=(X^{-1}(B))^c,\qquad X^{-1}\Bigl(\bigcup_iB_i\Bigr)=\bigcup_iX^{-1}(B_i),\qquad X^{-1}\Bigl(\bigcap_iB_i\Bigr)=\bigcap_iX^{-1}(B_i),
\]
对任意指标集成立。例如第一式：$\omega\in X^{-1}(B^c)\iff X(\omega)\notin B\iff\omega\notin X^{-1}(B)$。

\begin{definition}[可测映射，随机变量]\label{def:random-variables:measurable}
设 $(\Omega,\cF)$ 与 $(S,\cS)$ 是可测空间。映射 $X:\Omega\to S$ 若满足
\[
X^{-1}(B)\in\cF\qquad\text{对一切}\ B\in\cS,
\]
则称 $X$ 是从 $(\Omega,\cF)$ 到 $(S,\cS)$ 的\term{可测映射}（measurable map），也说 $X$ 是 $\cF/\cS$-可测的，或者在不致混淆时说 $X$ 是 \term{$\cF$-可测的}，记作 $X\in\cF$。当 $(\Omega,\cF,P)$ 是概率空间时，可测映射也称为\term{随机元}（random element）。特别地：
\begin{enumerate}
\item $(S,\cS)=(\R,\cR)$ 时，$X$ 称为\term{随机变量}（random variable）；
\item $(S,\cS)=(\R^d,\cR^d)$ 时，$X$ 称为 $d$ 维\term{随机向量}（random vector）；
\item $(\Omega,\cF)=(\R^m,\cR^m)$，$(S,\cS)=(\R^d,\cR^d)$ 时，$X$ 称为 \term{Borel 可测函数}（Borel measurable function），简称 Borel 函数。
\end{enumerate}
\end{definition}

请比较拓扑空间之间连续映射的定义：连续映射是开集的原像为开集的映射，可测映射是可测集的原像为可测集的映射。两个定义形式完全相同，只是把“开集的全体”换成了“$\sigma$-代数”。这个类比在下面很多地方都有帮助。

\begin{remark}[随机变量的概念与概率测度无关]
定义中根本没有出现 $P$。同一个函数 $X$，在同一个可测空间上配以不同的概率测度，就有不同的“随机规律”。例如在 $([0,1],\cR\cap[0,1])$ 上令 $X(\omega)=\omega$。若 $P$ 是 $[0,1]$ 上的 Lebesgue 测度，则 $P(X\leqslant x)=x$（$0\leqslant x\leqslant1$），$X$ 是“均匀分布的”；若 $P=\delta_{1/2}$，则 $X=1/2$ 几乎必然成立，$X$ 是“常数”。反过来，可测性依赖于 $\sigma$-代数：一个函数可以对 $\cF$ 可测，而对更小的 $\cG\subseteq\cF$ 不可测。以后讨论条件期望与鞅时，我们会在同一个 $\Omega$ 上同时使用多个 $\sigma$-代数，那时“对哪个 $\sigma$-代数可测”就成为核心问题。
\end{remark}

\begin{example}\label{ex:random-variables:basic}
\begin{enumerate}
\item 示性函数 $\one_A$ 是随机变量当且仅当 $A\in\cF$。事实上对 $B\in\cR$，$\{\one_A\in B\}$ 只可能是 $\varnothing$、$A$、$A^c$、$\Omega$ 之一（取决于 $0,1$ 是否属于 $B$），而 $\{\one_A\in\{1\}\}=A$。
\item 若 $\cF=2^\Omega$（例如 $\Omega$ 可数的离散概率空间），则 $\Omega$ 上的任何实值函数都是随机变量。
\item 若 $\cF=\{\varnothing,\Omega\}$，则只有常数是随机变量。设 $X$ 可测，则每个 $\{X\leqslant x\}$ 是 $\varnothing$ 或 $\Omega$。令 $c=\inf\{x:\{X\leqslant x\}=\Omega\}$。由于 $\Omega$ 非空，任取 $\omega_0$，$x<X(\omega_0)$ 时 $\{X\leqslant x\}\neq\Omega$，所以 $c\geqslant X(\omega_0)>-\infty$；而 $x=X(\omega_0)$ 时 $\{X\leqslant x\}\neq\varnothing$，从而等于 $\Omega$，所以 $c\leqslant X(\omega_0)<\infty$。于是 $c=X(\omega_0)$ 对每个 $\omega_0$ 成立，$X\equiv c$。
\end{enumerate}
\end{example}

\begin{workedexample}\label{we:random-variables:coin}
在无穷次抛硬币的空间 $\Omega=\{0,1\}^\N$，$\cF=\sigma(\{A_n\})$，$A_n=\{\omega_n=1\}$ 上（\egref{ex:set-systems:coin-tossing}），证明下列函数是（广义实值）随机变量：
\begin{enumerate}
\item $\xi_n(\omega)=\omega_n$（第 $n$ 次的结果）与 $S_n=\xi_1+\cdots+\xi_n$；
\item $T(\omega)=\inf\{n:\omega_n=1\}$（第一次出现正面的时刻，约定 $\inf\varnothing=\infty$）；
\item $L_n(\omega)$ = 前 $n$ 次中最长的连续正面的长度。
\end{enumerate}
\end{workedexample}
\begin{solution}
我们只用定义：对每个 $x$ 验证 $\{X\leqslant x\}\in\cF$，或者直接验证 $\{X\in B\}\in\cF$。（下面推论~\ref{cor:random-variables:criterion} 说明只验证 $\{X\leqslant x\}$ 就够了。）

(1) $\xi_n=\one_{A_n}$，对 Borel 集 $B$，$\{\xi_n\in B\}$ 是 $\varnothing,A_n,A_n^c,\Omega$ 之一，都属于 $\cF$。$S_n$ 只取 $0,1,\dots,n$ 这有限个值，$\{S_n=k\}$ 是前 $n$ 个坐标中恰有 $k$ 个 $1$ 的那些柱集之并，属于 $\cF_n\subseteq\cF$；于是 $\{S_n\in B\}=\bigcup_{k\in B\cap\{0,\dots,n\}}\{S_n=k\}\in\cF$。

(2) $T$ 取值于 $\{1,2,\dots\}\cup\{\infty\}$。对正整数 $n$，$T\leqslant n$ 当且仅当前 $n$ 次中至少有一次正面，所以
\[
\{T\leqslant n\}=A_1\cup\cdots\cup A_n\in\cF .
\]
对实数 $x$，$\{T\leqslant x\}=\{T\leqslant\lfloor x\rfloor\}$（$x<1$ 时为空集），所以 $\{T\leqslant x\}\in\cF$ 对一切实数 $x$ 成立，由下面的命题~\ref{prop:random-variables:extended}，$T$ 是广义实值随机变量。$\{T=\infty\}=\bigcap_nA_n^c$ 是“永远不出现正面”。

(3) $L_n(\omega)$ 只依赖于 $\omega_1,\dots,\omega_n$，只取 $0,\dots,n$ 这些值，$\{L_n=j\}$ 是有限个柱集的并，属于 $\cF_n$。与 (1) 同理，$L_n$ 是随机变量。

这个例题说明了一个一般的原则：\emph{只依赖于前 $n$ 次结果的量，都是 $\cF_n$-可测的随机变量}。这是因为 $\cF_n$ 恰好由前 $n$ 个坐标决定的柱集组成。至于依赖于全部无穷多次结果的量，比如 $\limsup_nS_n/n$，需要用下面的定理~\ref{thm:random-variables:limits}。
\end{solution}

\begin{theorem}[在生成元上验证可测性]\label{thm:random-variables:generators}
设 $X:\Omega\to S$，$\cC\subseteq2^S$，$\sigma(\cC)=\cS$。若 $X^{-1}(C)\in\cF$ 对一切 $C\in\cC$ 成立，则 $X$ 是 $\cF/\cS$-可测的。
\end{theorem}
\begin{proof}
令 $\cG=\{B\subseteq S: X^{-1}(B)\in\cF\}$。由原像与集合运算可交换：$X^{-1}(S)=\Omega\in\cF$，故 $S\in\cG$；若 $B\in\cG$，则 $X^{-1}(B^c)=(X^{-1}(B))^c\in\cF$，故 $B^c\in\cG$；若 $B_n\in\cG$，则 $X^{-1}(\bigcup_nB_n)=\bigcup_nX^{-1}(B_n)\in\cF$，故 $\bigcup_nB_n\in\cG$。所以 $\cG$ 是 $\sigma$-代数。由假设 $\cC\subseteq\cG$，由最小性 $\cS=\sigma(\cC)\subseteq\cG$，这正是可测性。
\end{proof}

结合命题~\ref{prop:set-systems:borel-generators}，立即得到：

\begin{corollary}\label{cor:random-variables:criterion}
$X:\Omega\to\R$ 是随机变量，当且仅当 $\{X\leqslant x\}\in\cF$ 对一切 $x\in\R$ 成立，也当且仅当 $\{X\leqslant x\}\in\cF$ 对一切 $x\in\Q$ 成立。同样，也可以用 $\{X<x\}$、$\{X>x\}$ 或 $\{X\geqslant x\}$ 代替 $\{X\leqslant x\}$。
\end{corollary}
\begin{proof}
由命题~\ref{prop:set-systems:borel-generators}，$\{(-\infty,x]:x\in\Q\}$ 生成 $\cR$，于是由定理~\ref{thm:random-variables:generators} 得第二个充分性；第一个充分性更弱；必要性是显然的。对其他三种形式：$\{X<x\}=\bigcup_n\{X\leqslant x-1/n\}$，$\{X\leqslant x\}=\bigcap_n\{X<x+1/n\}$，所以 $\{X\leqslant x\}\in\cF\ \forall x$ 与 $\{X<x\}\in\cF\ \forall x$ 等价；取补集得到另外两种。
\end{proof}

这个定理是“好集原理”的又一次应用：使 $X^{-1}(B)$ 可测的集合 $B$ 的全体总是 $\sigma$-代数，所以只需要对生成元验证。它的意义在于，$\cS$ 中的集合一般多得无法一一验证，而生成元往往很简单。

\begin{theorem}[复合与连续映射]\label{thm:random-variables:composition}
\begin{enumerate}
\item 若 $X:(\Omega,\cF)\to(S,\cS)$ 与 $f:(S,\cS)\to(T,\cT)$ 都可测，则 $f\circ X:(\Omega,\cF)\to(T,\cT)$ 可测。
\item 设 $S,T$ 是度量空间（或拓扑空间），$f:S\to T$ 连续，则 $f$ 是 $\cB(S)/\cB(T)$-可测的。
\item 单调函数 $f:\R\to\R$ 是 Borel 可测的。
\end{enumerate}
\end{theorem}
\begin{proof}
(1) 对 $B\in\cT$，$(f\circ X)^{-1}(B)=X^{-1}(f^{-1}(B))$，而 $f^{-1}(B)\in\cS$，故 $X^{-1}(f^{-1}(B))\in\cF$。

(2) 开集的原像是开集，从而属于 $\cB(S)$；开集生成 $\cB(T)$，由定理~\ref{thm:random-variables:generators} 即得。

(3) 设 $f$ 单调不减（单调不增时考虑 $-f$，并注意 $x\mapsto-x$ 连续）。对任意 $c$，集合 $I=\{x:f(x)\leqslant c\}$ 具有性质：若 $x\in I$ 且 $y<x$，则 $f(y)\leqslant f(x)\leqslant c$，$y\in I$。这样的集合只能是 $\varnothing$、$\R$、$(-\infty,a)$ 或 $(-\infty,a]$（$a=\sup I$），都是 Borel 集。由推论~\ref{cor:random-variables:criterion} 即得。
\end{proof}

\begin{remark}
(2) 的证明是简单的，但是其重要性不言而喻：拓扑空间之间的连续映射一定是可测的。一般而言，连续性是非常容易验证的，所以我们现在有了一大类可测函数的例子。再加上下面的四则运算与极限运算，分析中遇到的函数几乎都是可测的。事实上，要构造一个不可测的函数，必须借助选择公理（第~\ref{lec:extension-of-measures} 讲作业 D）。所以在实际的问题中，可测性很少是真正的障碍，但我们仍然需要知道它为什么成立。
\end{remark}

\begin{workedexample}\label{we:random-variables:borel-functions}
证明下列函数 $\R\to\R$ 是 Borel 函数：
\begin{enumerate}
\item Dirichlet 函数 $\one_\Q$；
\item $f(x)=\sin(1/x)$（$x\neq0$），$f(0)=0$；
\item $\lfloor x\rfloor$ 与 $x-\lfloor x\rfloor$。
\end{enumerate}
\end{workedexample}
\begin{solution}
(1) $\Q$ 是 Borel 集（可数个单点之并），由\egref{ex:random-variables:basic}(1)，$\one_\Q$ 是 Borel 函数。它处处不连续，所以这个例子说明可测函数可以非常不规则。

(2) $f$ 在 $0$ 处不连续，不能直接用定理~\ref{thm:random-variables:composition}(2)。用推论~\ref{cor:random-variables:criterion}：对实数 $a$，
\[
\{f\leqslant a\}=\{x\neq0:\ \sin(1/x)\leqslant a\}\cup\bigl(\{0\}\cap\{0\leqslant a\}\bigr).
\]
第二部分是 $\{0\}$ 或空集。第一部分：$g(x)=\sin(1/x)$ 在开集 $U=\R\setminus\{0\}$ 上连续，所以 $\{x\in U:g(x)>a\}$ 是 $U$ 中的相对开集，即某个开集与 $U$ 的交，也是开集；第一部分是 $U$ 减去这个开集，属于 $\cR$。所以 $\{f\leqslant a\}\in\cR$。一般地，若把 $\R$ 分成两个 Borel 集，函数在每一块上分别可测，则拼起来的函数可测，见习题~\exref{ex:random-variables:A1}。

(3) $\lfloor x\rfloor$ 单调不减，由定理~\ref{thm:random-variables:composition}(3) 是 Borel 函数；$x-\lfloor x\rfloor$ 是两个 Borel 函数之差，由下面的推论~\ref{cor:random-variables:arithmetic} 是 Borel 函数。
\end{solution}

\begin{theorem}[随机向量与分量]\label{thm:random-variables:vector}
设 $X=(X_1,\dots,X_d):\Omega\to\R^d$。则 $X$ 是随机向量，当且仅当每个 $X_i$ 都是随机变量。
\end{theorem}
\begin{proof}
必要性：坐标投影 $p_i(x)=x_i$ 是连续映射 $\R^d\to\R$，由定理~\ref{thm:random-variables:composition}，$X_i=p_i\circ X$ 可测。充分性：对 $x\in\R^d$，
\[
\{X\in(-\infty,x_1]\times\cdots\times(-\infty,x_d]\}=\bigcap_{i=1}^d\{X_i\leqslant x_i\}\in\cF,
\]
而这样的集合族生成 $\cR^d$（习题~\exref{ex:set-systems:A5}），由定理~\ref{thm:random-variables:generators} 即得。
\end{proof}

\begin{corollary}\label{cor:random-variables:arithmetic}
设 $X_1,\dots,X_n$ 是随机变量，$f:(\R^n,\cR^n)\to(\R,\cR)$ 是 Borel 函数（例如连续函数），则 $f(X_1,\dots,X_n)$ 是随机变量。特别地，$X_1+X_2$，$X_1X_2$，$cX_1$，$\max\{X_1,X_2\}$，$\min\{X_1,X_2\}$，$X_1^+=\max\{X_1,0\}$，$X_1^-=\max\{-X_1,0\}$，$\abs{X_1}$ 都是随机变量；若 $X_2$ 处处不为 $0$，则 $X_1/X_2$ 是随机变量。
\end{corollary}
\begin{proof}
由定理~\ref{thm:random-variables:vector}，$(X_1,\dots,X_n)$ 是随机向量，由定理~\ref{thm:random-variables:composition}(1) 即得第一句。所列的函数 $(x,y)\mapsto x+y,\ xy,\ \max\{x,y\}$ 等在 $\R^2$ 上连续，由定理~\ref{thm:random-variables:composition}(2) 是 Borel 函数。对商：令 $g(y)=1/y$（$y\neq0$），$g(0)=0$。对 $c>0$，$\{g\leqslant c\}=(-\infty,0]\cup[1/c,\infty)$；对 $c=0$，$\{g\leqslant0\}=(-\infty,0]$；对 $c<0$，$\{g\leqslant c\}=[1/c,0)$。它们都是 Borel 集，所以 $g$ 是 Borel 函数，$X_1/X_2=X_1\cdot g(X_2)$。
\end{proof}

我们再换一个角度看 $X_1+X_2$ 的可测性。上面的证明把 $(X_1,X_2)$ 看成取值于 $\R^2$ 的随机向量，再复合上连续函数 $(x,y)\mapsto x+y$。其实也可以直接验证：对任意 $c\in\R$，
\[
\{X_1+X_2<c\}=\bigcup_{q\in\Q}\bigl(\{X_1<q\}\cap\{X_2<c-q\}\bigr).
\]
这因为：若 $X_1(\omega)+X_2(\omega)<c$，则 $X_1(\omega)<c-X_2(\omega)$，由有理数的稠密性，可以取有理数 $q$ 使 $X_1(\omega)<q<c-X_2(\omega)$，此时 $X_2(\omega)<c-q$；反之，若 $X_1(\omega)<q$ 且 $X_2(\omega)<c-q$，两式相加得 $X_1(\omega)+X_2(\omega)<c$。右边是可数个事件的并，所以属于 $\cF$，再由推论~\ref{cor:random-variables:criterion} 即得。这个直接的证明用到了有理数的稠密性，每换一个函数（比如 $X_1X_2$）就要重新想一个办法。前面那个经过 $\R^2$ 的证明则一劳永逸：只要 $f$ 是 $\R^2$ 上的 Borel 函数，$f(X_1,X_2)$ 就是随机变量。为了把这个想法用到一般的可测空间上，我们需要乘积空间上的 $\sigma$-代数。

\topic{乘积 \texorpdfstring{$\sigma$}{σ}-代数与截口}

\begin{definition}[乘积 $\sigma$-代数]\label{def:random-variables:product}
设 $(S,\cS)$ 与 $(T,\cT)$ 是可测空间。形如 $A\times B$（$A\in\cS$，$B\in\cT$）的集合称为 $S\times T$ 中的\term{可测矩形}（measurable rectangle）。由全体可测矩形生成的 $\sigma$-代数
\[
\cS\otimes\cT=\sigma\bigl(\{A\times B: A\in\cS,\ B\in\cT\}\bigr)
\]
称为\term{乘积 $\sigma$-代数}（product $\sigma$-algebra）。有限个可测空间的乘积 $\cS_1\otimes\cdots\otimes\cS_d$ 同样定义。
\end{definition}

两个可测矩形的交还是可测矩形：$(A\times B)\cap(A'\times B')=(A\cap A')\times(B\cap B')$，所以可测矩形的全体是一个 $\pi$-类。以后构造乘积测度时，我们会用它和唯一性定理。我们强调，$\cS\otimes\cT$ 中的集合一般不是矩形。比如说平面上的对角线 $\{(x,x):x\in\R\}$ 是闭集，属于 $\cR^2$，但它不是任何 $A\times B$ 的形式：若它等于 $A\times B$，则 $(0,0)$ 与 $(1,1)$ 都在其中，于是 $0\in A$，$1\in B$，从而 $(0,1)\in A\times B$，矛盾。

\begin{proposition}\label{prop:random-variables:Rd-product}
$\cR^d=\cR\otimes\cdots\otimes\cR$（$d$ 个）。
\end{proposition}
\begin{proof}
记右边为 $\cR^{\otimes d}$。一方面，每个矩形 $(a_1,b_1]\times\cdots\times(a_d,b_d]$ 都是可测矩形，而这样的矩形生成 $\cR^d$（习题~\exref{ex:set-systems:A5}），所以 $\cR^d\subseteq\cR^{\otimes d}$。另一方面，设 $A_1,\dots,A_d\in\cR$。坐标投影 $p_i(x)=x_i$ 是 $\R^d\to\R$ 的连续映射，由定理~\ref{thm:random-variables:composition}(2)，它是 $\cR^d/\cR$-可测的，所以
\[
A_1\times\cdots\times A_d=\bigcap_{i=1}^dp_i^{-1}(A_i)\in\cR^d .
\]
于是 $\cR^d$ 包含一切可测矩形，从而 $\cR^{\otimes d}\subseteq\cR^d$。
\end{proof}

\begin{theorem}[取值于乘积空间的映射]\label{thm:random-variables:product-map}
设 $X:\Omega\to S$，$Y:\Omega\to T$，$(X,Y):\Omega\to S\times T$ 表示 $\omega\mapsto(X(\omega),Y(\omega))$。则 $(X,Y)$ 是 $\cF/\cS\otimes\cT$-可测的，当且仅当 $X$ 是 $\cF/\cS$-可测的且 $Y$ 是 $\cF/\cT$-可测的。
\end{theorem}
\begin{proof}
必要性：投影 $\pi_S(s,t)=s$ 是 $\cS\otimes\cT/\cS$-可测的，因为 $\pi_S^{-1}(A)=A\times T$ 是可测矩形。于是 $X=\pi_S\circ(X,Y)$ 可测（定理~\ref{thm:random-variables:composition}(1)）。$Y$ 同理。

充分性：对可测矩形 $A\times B$，$\{(X,Y)\in A\times B\}=\{X\in A\}\cap\{Y\in B\}\in\cF$。可测矩形生成 $\cS\otimes\cT$，由定理~\ref{thm:random-variables:generators} 即得。
\end{proof}

由命题~\ref{prop:random-variables:Rd-product}，定理~\ref{thm:random-variables:vector} 是这个定理的特例。这个定理说明，乘积 $\sigma$-代数恰好是使“联合起来可测”与“分别可测”等价的那个 $\sigma$-代数：它是使两个投影都可测的最小 $\sigma$-代数（习题~\exref{ex:random-variables:A10}）。于是对任意 $\cS\otimes\cT$-可测的函数 $f:S\times T\to\R$，$f(X,Y)$ 都是随机变量。

下面讨论乘积空间中集合的截口。对 $E\subseteq S\times T$ 与 $s\in S$，称
\[
E_s=\{t\in T:(s,t)\in E\}
\]
为 $E$ 在 $s$ 处的\term{截口}（section）；类似地，对 $t\in T$，$E^t=\{s\in S:(s,t)\in E\}$。对函数 $f:S\times T\to\R$，令 $f_s(t)=f(s,t)$，$f^t(s)=f(s,t)$。几何上，$E_s$ 就是用竖直的直线 $\{s\}\times T$ 去截 $E$ 所得的集合。

\begin{theorem}[截口的可测性]\label{thm:random-variables:sections}
设 $E\in\cS\otimes\cT$，$f:S\times T\to\R$ 是 $\cS\otimes\cT$-可测的。则对每个 $s\in S$，$E_s\in\cT$，$f_s$ 是 $\cT$-可测的；对每个 $t\in T$，$E^t\in\cS$，$f^t$ 是 $\cS$-可测的。
\end{theorem}
\begin{proof}
固定 $s\in S$，令 $\iota_s:T\to S\times T$，$\iota_s(t)=(s,t)$。对可测矩形 $A\times B$，
\[
\iota_s^{-1}(A\times B)=\{t:(s,t)\in A\times B\}=\begin{cases}B,&s\in A,\\ \varnothing,&s\notin A,\end{cases}
\]
都属于 $\cT$。由定理~\ref{thm:random-variables:generators}，$\iota_s$ 是 $\cT/\cS\otimes\cT$-可测的。而按照定义，$E_s=\iota_s^{-1}(E)$，$f_s=f\circ\iota_s$，所以 $E_s\in\cT$，并且由定理~\ref{thm:random-variables:composition}(1)，$f_s$ 是 $\cT$-可测的。$E^t$ 与 $f^t$ 同理。
\end{proof}

特别地，取 $S=T=\R$，由命题~\ref{prop:random-variables:Rd-product}：\emph{平面上 Borel 集的截口都是直线上的 Borel 集，平面上 Borel 函数固定一个变量后是另一个变量的 Borel 函数。}这个事实在以后的 Fubini 定理中是第一步：要计算 $\int\bigl(\int f(s,t)\,dt\bigr)ds$，首先要知道里面的积分对每个 $s$ 有意义。

\begin{remark}[Lebesgue 可测集的截口]
这个结论对 Lebesgue 可测集并不成立。回忆第~\ref{lec:extension-of-measures} 讲，Lebesgue 可测集（外测度 $\lambda^*$ 的 Carathéodory 可测集）是比 Borel 集更大的一族集合，零外测度集的每个子集都是 Lebesgue 可测的；用二维外测度 $\lambda_2^*$ 同样可以定义平面上的 Lebesgue 可测集。设 $V\subseteq[0,1)$ 不是 Lebesgue 可测集（第~\ref{lec:extension-of-measures} 讲作业 D），令 $E=\{0\}\times V$。$E$ 包含于直线 $\{0\}\times\R$，后者的二维 Lebesgue 测度为零（第~\ref{lec:extension-of-measures} 讲作业 B6），所以 $E$ 的二维外测度为零，是平面上的 Lebesgue 可测集。可是它在 $s=0$ 处的截口 $E_0=V$ 不是 Lebesgue 可测的。问题出在：Lebesgue 可测集的全体不是由矩形生成的，它还把一切零测集的子集都放了进来，而零测集的子集可以非常古怪。这是概率论中宁愿使用 Borel $\sigma$-代数与乘积 $\sigma$-代数的原因之一。
\end{remark}

\topic{广义实值随机变量与极限运算}

一列随机变量的上确界、上极限可能等于 $\pm\infty$，比如 $\sup_nX_n$。为了让它们也成为随机变量，我们把取值范围扩充为 $[-\infty,\infty]$。记 $\overline\R=[-\infty,\infty]$，令
\[
\overline\cR=\{B\subseteq\overline\R:\ B\cap\R\in\cR\}.
\]
容易验证 $\overline\cR$ 是 $\overline\R$ 上的 $\sigma$-代数（$(B^c)\cap\R=\R\setminus(B\cap\R)$，$(\bigcup_nB_n)\cap\R=\bigcup_n(B_n\cap\R)$），它包含 $\cR$ 以及 $\{\infty\}$、$\{-\infty\}$。（可以证明它就是 $\overline\R$ 在通常拓扑下的 Borel $\sigma$-代数，但我们用不到这一点。）取值于 $\overline\R$ 的 $\cF/\overline\cR$-可测函数称为\term{广义实值随机变量}（extended real-valued random variable）。以后在不致混淆时也简称随机变量。

\begin{proposition}\label{prop:random-variables:extended}
$X:\Omega\to\overline\R$ 是 $\cF/\overline\cR$-可测的，当且仅当 $\{X\leqslant x\}\in\cF$ 对一切 $x\in\R$ 成立。
\end{proposition}
\begin{proof}
必要性：$[-\infty,x]\in\overline\cR$。充分性：令 $\cG=\{B\subseteq\overline\R: X^{-1}(B)\in\cF\}$，与定理~\ref{thm:random-variables:generators} 的证明相同，$\cG$ 是 $\sigma$-代数。由假设，$[-\infty,x]\in\cG$。于是 $\{-\infty\}=\bigcap_n[-\infty,-n]\in\cG$，$\{\infty\}=\bigcap_n(\overline\R\setminus[-\infty,n])\in\cG$，$(a,b]=[-\infty,b]\setminus[-\infty,a]\in\cG$。$\cG\cap\cR=\{B\in\cR: B\in\cG\}$ 是 $\R$ 上的 $\sigma$-代数（因为 $\R=\overline\R\setminus\{\pm\infty\}\in\cG$），它含有一切 $(a,b]$，所以含有 $\cR$。最后，任取 $B\in\overline\cR$，$B=(B\cap\R)\cup(B\cap\{-\infty\})\cup(B\cap\{\infty\})$ 是 $\cG$ 中至多三个集合的并，故 $B\in\cG$。
\end{proof}

\begin{theorem}[上下确界与上下极限]\label{thm:random-variables:limits}
设 $X_1,X_2,\dots$ 是（广义实值）随机变量。则 $\inf_nX_n$，$\sup_nX_n$，$\liminf_nX_n$，$\limsup_nX_n$ 都是广义实值随机变量。进一步，
\[
\Omega_0=\{\omega:\ \lim_nX_n(\omega)\ \text{存在且有限}\}\in\cF,
\]
并且若 $\Omega_0=\Omega$，则 $\lim_nX_n$ 是随机变量。
\end{theorem}
\begin{proof}
对 $x\in\R$，$\sup_nX_n(\omega)\leqslant x$ 当且仅当对一切 $n$ 有 $X_n(\omega)\leqslant x$，所以 $\{\sup_nX_n\leqslant x\}=\bigcap_n\{X_n\leqslant x\}\in\cF$，由命题~\ref{prop:random-variables:extended}，$\sup_nX_n$ 可测。$\inf_nX_n(\omega)<x$ 当且仅当某个 $X_n(\omega)<x$，所以 $\{\inf_nX_n<x\}=\bigcup_n\{X_n<x\}\in\cF$，从而 $\{\inf_nX_n\leqslant x\}=\bigcap_m\{\inf_nX_n<x+1/m\}\in\cF$。

由于 $\liminf_nX_n=\sup_n\inf_{m\geqslant n}X_m$，$\limsup_nX_n=\inf_n\sup_{m\geqslant n}X_m$，对每个 $n$，$Y_n=\inf_{m\geqslant n}X_m$ 可测，再取上确界仍可测；$\limsup$ 同理。

$\lim_nX_n(\omega)$ 存在且有限，当且仅当 $U(\omega)=\limsup_nX_n(\omega)$ 与 $L(\omega)=\liminf_nX_n(\omega)$ 相等且有限。总有 $L\leqslant U$。由于 $L(\omega)<U(\omega)$ 当且仅当存在有理数 $q$ 使 $L(\omega)<q<U(\omega)$，
\[
\{L<U\}=\bigcup_{q\in\Q}\bigl(\{L<q\}\cap\{U>q\}\bigr)\in\cF .
\]
又 $\{U<\infty\}=\bigcup_n\{U\leqslant n\}\in\cF$，$\{L>-\infty\}=\bigcup_n\{L\geqslant-n\}\in\cF$（这里 $\{L\geqslant-n\}=\{L<-n\}^c$，而 $\{L<-n\}=\bigcup_m\{L\leqslant-n-1/m\}$）。所以
\[
\Omega_0=\{U<\infty\}\cap\{L>-\infty\}\cap\{L<U\}^c\in\cF .
\]
若 $\Omega_0=\Omega$，则 $\lim_nX_n=\limsup_nX_n$ 取有限值，并且可测。
\end{proof}

\begin{remark}
这个定理说明，可测函数的全体对逐点极限是封闭的。与连续函数对比：连续函数列的逐点极限不一定连续（$[0,1]$ 上的 $x^n$ 的极限在 $x=1$ 处跳跃），可测函数列的逐点极限却总是可测的。可测性之所以经得起极限，归根到底是因为 $\sigma$-代数对可数次运算封闭，而 $\{\sup_nX_n\leqslant x\}=\bigcap_n\{X_n\leqslant x\}$ 只用到可数交。这是积分理论要建立在可测函数而不是连续函数上的一个重要原因：分析中大量的对象是通过取极限得到的，我们需要一个对极限封闭的函数类。
\end{remark}

\begin{definition}[几乎必然收敛]\label{def:random-variables:as}
设 $(\Omega,\cF,P)$ 是概率空间，$X_n,X$ 是随机变量。若 $P(\{\omega:X_n(\omega)\to X(\omega)\})=1$，称 $X_n$ \term{几乎必然收敛}（converges almost surely）于 $X$，记作 $X_n\to X$ a.s.
\end{definition}

由定理~\ref{thm:random-variables:limits} 的证明，$\{X_n\to X\}=\bigcap_{k}\bigcup_N\bigcap_{n\geqslant N}\{\abs{X_n-X}<1/k\}\in\cF$，所以这个定义中的概率有意义。当极限只在一个零测集之外存在时，常令 $X_\infty=\limsup_nX_n$，它是广义实值随机变量，a.s. 等于极限。

\topic{不可数个随机变量的上确界}

定理~\ref{thm:random-variables:limits} 中取上确界的是可数个随机变量，证明中用到的是 $\{\sup_nX_n\leqslant x\}=\bigcap_n\{X_n\leqslant x\}$ 这个可数交。在概率论中我们经常遇到不可数多个随机变量，比如说随时间连续变化的随机过程 $\{X_t:t\in[0,1]\}$，我们关心它在整个时间段上的最大值 $\sup_{t\in[0,1]}X_t$。这时
\[
\Bigl\{\sup_{t\in[0,1]}X_t\leqslant x\Bigr\}=\bigcap_{t\in[0,1]}\{X_t\leqslant x\}
\]
是不可数个事件的交，而 $\sigma$-代数对不可数交并不封闭。于是一个自然的问题是：不可数个随机变量的上确界、下确界是不是随机变量？一般来说，答案是否定的。

\begin{example}[不可数个随机变量的上确界可以不可测]
取 $(\Omega,\cF)=(\R,\cR)$。由第~\ref{lec:set-systems} 讲作业 F，存在不是 Borel 集的集合 $V\subseteq\R$。对 $t\in V$ 令 $X_t=\one_{\{t\}}$。单点集是闭集，所以每个 $X_t$ 都是随机变量。可是对每个 $\omega$，
\[
\sup_{t\in V}X_t(\omega)=\begin{cases}1,&\omega\in V,\\0,&\omega\notin V,\end{cases}
\]
即 $\sup_{t\in V}X_t=\one_V$，由\egref{ex:random-variables:basic} 的 (1)，它不是随机变量。同样，$\inf_{t\in V}(1-X_t)=\one_{V^c}$ 也不是随机变量。这个例子里，每个 $X_t$ 都再简单不过，毛病完全出在“不可数个”上。
\end{example}

所以要保证不可数个随机变量的上确界可测，必须对 $t\mapsto X_t(\omega)$ 的正则性有所要求。第~\ref{lec:measures} 讲 $C[0,1]$ 的例子中我们已经见过这种想法：对连续函数，$\sup_{t\in[0,1]}f(t)=\sup_{t\in\Q\cap[0,1]}f(t)$，不可数个时刻可以用可数个有理时刻代替。下面的命题说，只要右连续就够了。

\begin{proposition}[右连续时化为可数个]\label{prop:random-variables:sup-right-continuous}
设 $T\subseteq\R$，$D\subseteq T$ 是可数集，并且对每个 $t\in T$，存在 $d_n\in D$ 使 $d_n\geqslant t$ 且 $d_n\to t$。（例如 $T=\R$，$D=\Q$；或者 $T=[a,b]$，$D=(\Q\cap[a,b])\cup\{b\}$。）设 $\{X_t:t\in T\}$ 是一族随机变量，并且对每个 $\omega$，函数 $t\mapsto X_t(\omega)$ 在 $T$ 上右连续，即 $s\in T$，$s\geqslant t$，$s\to t$ 时 $X_s(\omega)\to X_t(\omega)$。则
\[
\sup_{t\in T}X_t=\sup_{t\in D}X_t,\qquad \inf_{t\in T}X_t=\inf_{t\in D}X_t,
\]
从而它们都是广义实值随机变量。
\end{proposition}
\begin{proof}
固定 $\omega$。由于 $D\subseteq T$，$\sup_{t\in D}X_t(\omega)\leqslant\sup_{t\in T}X_t(\omega)$。反之，任取 $t\in T$，按假设取 $d_n\in D$，$d_n\geqslant t$，$d_n\to t$，由右连续性，
\[
X_t(\omega)=\lim_nX_{d_n}(\omega)\leqslant\sup_{s\in D}X_s(\omega).
\]
对 $t\in T$ 取上确界得 $\sup_{t\in T}X_t(\omega)\leqslant\sup_{s\in D}X_s(\omega)$。下确界同理。$D$ 是可数集，由定理~\ref{thm:random-variables:limits}，$\sup_{t\in D}X_t$ 与 $\inf_{t\in D}X_t$ 是广义实值随机变量。
\end{proof}

我们来看一个统计中的重要例子。设 $X_1,\dots,X_n$ 是概率空间 $(\Omega,\cF,P)$ 上的随机变量，我们把它们想成对某个量的 $n$ 次观测，比如 $n$ 个人的身高。令
\[
F_n(x)=F_n(x,\omega)=\frac1n\sum_{i=1}^n\one_{\{X_i(\omega)\leqslant x\}},\qquad x\in\R,
\]
即 $X_1(\omega),\dots,X_n(\omega)$ 中不超过 $x$ 的个数所占的比例。它称为\term{经验分布函数}（empirical distribution function）。对固定的 $x$，$F_n(x,\cdot)$ 是 $n$ 个示性函数的平均，是随机变量；对固定的 $\omega$，$x\mapsto F_n(x,\omega)$ 是一个分布函数，它对应的分布在每个观测值 $X_i(\omega)$ 处放上质量 $1/n$，是一个在 $X_i(\omega)$ 处跳跃的阶梯函数，单调不减且右连续。

在统计中，$X_i$ 的分布函数 $F$ 往往是未知的，而 $F_n$ 是可以从数据算出来的。自然要问 $F_n$ 与 $F$ 相差多少。衡量两个函数的差别，最直接的办法是看最大的偏差
\[
D_n=\sup_{x\in\R}\abs{F_n(x)-F(x)}.
\]
这里对不可数个 $x$ 取了上确界，所以 $D_n$ 是不是随机变量，需要验证。如果它不是随机变量，那么“$D_n\to0$ 的概率是多少”这个问题就没有意义。

\begin{proposition}\label{prop:random-variables:empirical}
设 $X_1,\dots,X_n$ 是随机变量，$F$ 是任意一个分布函数。则 $D_n=\sup_{x\in\R}\abs{F_n(x)-F(x)}$ 是取值于 $[0,1]$ 的随机变量。对一列随机变量 $X_1,X_2,\dots$，事件 $\{D_n\to0\}$ 属于 $\cF$。
\end{proposition}
\begin{proof}
令 $Y_x(\omega)=\abs{F_n(x,\omega)-F(x)}$。对每个 $x$，$F_n(x,\cdot)$ 是随机变量，$F(x)$ 是常数，由推论~\ref{cor:random-variables:arithmetic}，$Y_x$ 是随机变量。对每个 $\omega$，$x\mapsto F_n(x,\omega)$ 与 $x\mapsto F(x)$ 都右连续，二者之差右连续，再复合上连续函数 $\abs{\cdot}$ 仍右连续。由命题~\ref{prop:random-variables:sup-right-continuous}（取 $T=\R$，$D=\Q$），$D_n=\sup_{x\in\Q}Y_x$ 是广义实值随机变量。由于 $F_n$ 与 $F$ 都取值于 $[0,1]$，$0\leqslant Y_x\leqslant1$，所以 $D_n$ 取值于 $[0,1]$，是随机变量。

最后，$D_n\geqslant0$，所以 $D_n(\omega)\to0$ 当且仅当 $\limsup_nD_n(\omega)\leqslant0$，而由定理~\ref{thm:random-variables:limits}，$\limsup_nD_n$ 是广义实值随机变量，故 $\{D_n\to0\}=\{\limsup_nD_n\leqslant0\}\in\cF$。
\end{proof}

\begin{remark}[Glivenko--Cantelli 定理]
以后学了独立性与强大数定律，我们会证明：若 $X_1,X_2,\dots$ 独立同分布，分布函数为 $F$，则 $D_n\to0$ a.s.。这称为 \term{Glivenko--Cantelli 定理}，在统计中有时被称为统计学的基本定理：观测次数足够多时，经验分布函数在整条实轴上一致地逼近真实的分布函数。对每个固定的 $x$，$F_n(x)\to F(x)$ a.s. 只是强大数定律；难处在于 $x$ 有不可数多个，而每个 $x$ 都有一个概率为零的例外集，不可数个零测集的并可以不是零测集。证明时要用 $F_n$ 与 $F$ 的单调性，把不可数个 $x$ 归结为有限个点。这与上面把 $\sup_{x\in\R}$ 化成 $\sup_{x\in\Q}$ 是同一种想法。
\end{remark}

\topic{由随机变量生成的 \texorpdfstring{$\sigma$}{σ}-代数}

\begin{definition}[生成的 $\sigma$-代数]\label{def:random-variables:sigma-X}
设 $X:\Omega\to S$，$(S,\cS)$ 是可测空间。称
\[
\sigma(X)=X^{-1}(\cS)=\{\{X\in B\}: B\in\cS\}
\]
为由 $X$ \term{生成的 $\sigma$-代数}（$\sigma$-algebra generated by $X$）。对一族映射 $X_i:\Omega\to S_i$（$i\in I$），令 $\sigma(X_i:i\in I)=\sigma\bigl(\bigcup_{i\in I}\sigma(X_i)\bigr)$。
\end{definition}

由原像与集合运算可交换，$\sigma(X)$ 确实是 $\sigma$-代数（$\Omega=X^{-1}(S)$，$\{X\in B\}^c=\{X\in B^c\}$，$\bigcup_n\{X\in B_n\}=\{X\in\bigcup_nB_n\}$）。$X$ 是 $\cG/\cS$-可测的，当且仅当 $\sigma(X)\subseteq\cG$；所以 $\sigma(X)$ 是使 $X$ 可测的\uline{最小}的 $\sigma$-代数。直观上，我们总是把 $\sigma(X)$ 理解为“观察到 $X$ 的值所获得的信息”：一个事件属于 $\sigma(X)$，意思是只要看到 $X$ 的值，就能判断它是否发生。比如说，把硬币抛两次，$X$ 是第一次的结果（正面记为 $1$，反面记为 $0$），那么 $\sigma(X)=\{\varnothing,\{\text{正正},\text{正反}\},\{\text{反正},\text{反反}\},\Omega\}$。“第一次是正面”在其中，“两次结果相同”不在其中，因为只看第一次的结果无法判断两次是否相同。这与第~\ref{lec:set-systems} 讲中划分生成的 $\sigma$-代数是同一回事。

\begin{proposition}\label{prop:random-variables:sigma-generators}
若 $\cC$ 生成 $\cS$，则 $\sigma(X)=\sigma(\{X\in C\}:C\in\cC)$。特别地，对随机变量 $X$，$\sigma(X)=\sigma(\{X\leqslant x\}:x\in\R)$。
\end{proposition}
\begin{proof}
记右边为 $\cG$。由于 $\{X\in C\}\in\sigma(X)$，$\cG\subseteq\sigma(X)$。反之，$X^{-1}(C)\in\cG$ 对 $C\in\cC$ 成立，由定理~\ref{thm:random-variables:generators}，$X$ 是 $\cG/\cS$-可测的，所以 $\sigma(X)\subseteq\cG$。
\end{proof}

\begin{example}\label{ex:random-variables:sigma-examples}
\begin{enumerate}
\item $\sigma(\one_A)=\{\varnothing,A,A^c,\Omega\}$。
\item 若 $X$ 只取有限个不同的值 $x_1,\dots,x_k$，则 $\sigma(X)$ 是由划分 $\{X=x_1\},\dots,\{X=x_k\}$ 生成的 $\sigma$-代数（\egref{ex:set-systems:partition}）：$\{X\in B\}=\bigcup_{i:x_i\in B}\{X=x_i\}$。
\item 在 $\Omega=[-1,1]$ 上令 $X(\omega)=\omega^2$。则 $\{X\in B\}$ 总是关于原点对称的集合，所以 $\sigma(X)$ 中不含 $[0,1]$。直观上，知道 $\omega^2$ 不能判断 $\omega$ 的正负。
\end{enumerate}
\end{example}

\begin{workedexample}\label{we:random-variables:sigma-sum}
把公平硬币抛两次，$\Omega=\{00,01,10,11\}$，$\cF=2^\Omega$，$\omega=\omega_1\omega_2$。令 $X=\omega_1+\omega_2$（正面的次数），$Y=\omega_1$（第一次的结果），$Z=\one_{\{\omega_1\neq\omega_2\}}$（两次结果不同）。写出 $\sigma(X)$，并判断 $Y$、$Z$ 是否是 $\sigma(X)$-可测的。
\end{workedexample}
\begin{solution}
$X$ 取值 $0,1,2$，由\egref{ex:random-variables:sigma-examples}(2)，$\sigma(X)$ 由划分
\[
\{X=0\}=\{00\},\qquad\{X=1\}=\{01,10\},\qquad\{X=2\}=\{11\}
\]
生成，是这三块的并的全体，共 $2^3=8$ 个元素。

$Y$ 不是 $\sigma(X)$-可测的：$\{Y=1\}=\{10,11\}$ 不是三块中某几块的并（它把 $\{01,10\}$ 这一块切开了）。直观上，知道一共有一次正面，不能知道是第一次还是第二次。

$Z$ 是 $\sigma(X)$-可测的：$\{Z=1\}=\{01,10\}=\{X=1\}$。事实上 $Z=f(X)$，其中 $f(1)=1$，$f(0)=f(2)=0$。这正是下面 Doob--Dynkin 引理（定理~\ref{thm:random-variables:doob-dynkin}）的内容：$\sigma(X)$-可测的随机变量恰好是 $X$ 的函数。判断方法也很简单：$Z$ 在每一块上取常数，而 $Y$ 在 $\{01,10\}$ 这一块上取两个不同的值。
\end{solution}

\topic{简单函数逼近}

\begin{definition}[简单函数]\label{def:random-variables:simple}
形如 $\varphi=\sum_{i=1}^nc_i\one_{A_i}$（$c_i\in\R$，$A_i\in\cF$）的函数称为\term{简单函数}（simple function）。
\end{definition}

简单函数是随机变量（推论~\ref{cor:random-variables:arithmetic}），只取有限个值。反之，只取有限个值 $y_1,\dots,y_m$ 的随机变量 $Y$ 可以写成 $\sum_jy_j\one_{\{Y=y_j\}}$，其中 $\{Y=y_j\}$ 两两不交。

\begin{theorem}[简单函数逼近]\label{thm:random-variables:simple-approx}
设 $X:\Omega\to[0,\infty]$ 是 $\cF$-可测的。令
\[
X_n=\sum_{k=0}^{n2^n-1}\frac{k}{2^n}\one_{\{k2^{-n}\leqslant X<(k+1)2^{-n}\}}+n\one_{\{X\geqslant n\}}.
\]
则每个 $X_n$ 是非负简单函数，$0\leqslant X_1\leqslant X_2\leqslant\cdots$，且对每个 $\omega$，$X_n(\omega)\uparrow X(\omega)$。在 $\{X\leqslant M\}$ 上，$n\geqslant M$ 时 $0\leqslant X-X_n\leqslant2^{-n}$，所以收敛在 $\{X\leqslant M\}$ 上是一致的。
\end{theorem}
\begin{proof}
各个集合 $\{k2^{-n}\leqslant X<(k+1)2^{-n}\}=\{X\geqslant k2^{-n}\}\cap\{X<(k+1)2^{-n}\}$ 与 $\{X\geqslant n\}$ 都属于 $\cF$，故 $X_n$ 是简单函数。换一种写法：$X_n=\min\{n,\ 2^{-n}\lfloor2^nX\rfloor\}$（约定 $\lfloor\infty\rfloor=\infty$）。事实上若 $X(\omega)\geqslant n$，两边都是 $n$；若 $k2^{-n}\leqslant X(\omega)<(k+1)2^{-n}$，$k\leqslant n2^n-1$，则 $\lfloor2^nX(\omega)\rfloor=k$，两边都是 $k2^{-n}$。

\emph{单调性。}对实数 $t\geqslant0$，$2\lfloor t\rfloor\leqslant\lfloor2t\rfloor$（因为 $2\lfloor t\rfloor$ 是不超过 $2t$ 的整数），于是 $2^{-n}\lfloor2^nX\rfloor=2^{-(n+1)}\cdot2\lfloor2^nX\rfloor\leqslant2^{-(n+1)}\lfloor2^{n+1}X\rfloor$。又 $n\leqslant n+1$，取最小值保持不等式，故 $X_n\leqslant X_{n+1}$。

\emph{收敛。}若 $X(\omega)=\infty$，则 $X_n(\omega)=n\to\infty$。若 $X(\omega)<\infty$，当 $n\geqslant X(\omega)$ 时，$X_n(\omega)=2^{-n}\lfloor2^nX(\omega)\rfloor$，而 $2^nX(\omega)-1<\lfloor2^nX(\omega)\rfloor\leqslant2^nX(\omega)$，所以 $X(\omega)-2^{-n}<X_n(\omega)\leqslant X(\omega)$。这同时给出了 $\{X\leqslant M\}$ 上的一致估计。
\end{proof}

\begin{remark}[分割值域]
定理中的构造值得仔细体会。Riemann 积分是把\emph{定义域}分割成小区间，在每个小区间上用常数代替函数；这里是把\emph{值域}分割成长度为 $2^{-n}$ 的小段 $[k2^{-n},(k+1)2^{-n})$，看 $X$ 落在哪一段，用该段的左端点代替 $X$。后一种做法对 $\Omega$ 没有任何要求，$\Omega$ 可以是任何抽象的集合，只要 $\{k2^{-n}\leqslant X<(k+1)2^{-n}\}$ 是可测集就可以了。这就是 Lebesgue 积分的出发点。Lebesgue 本人打过一个比方：数一堆硬币的钱数，Riemann 的做法是按拿到手的顺序一枚一枚地加，Lebesgue 的做法是先把面值相同的硬币归成一堆，数出每堆有几枚，再乘以面值相加。
\end{remark}

对一般的实值随机变量 $X$，写 $X=X^+-X^-$，分别对 $X^+,X^-$ 用定理~\ref{thm:random-variables:simple-approx}，得到简单函数 $\varphi_n=X_n^+-X_n^-\to X$，并且 $\abs{\varphi_n}\leqslant\abs X$。

\topic{函数形式的单调类定理}

我们现在要把 $\pi$-$\lambda$ 定理从集合搬到函数上。这件事的动机是：很多时候我们要证明的不是“某个性质对 $\sigma(\cP)$ 中的每个集合成立”，而是“某个性质对每个有界的 $\sigma(\cP)$-可测函数成立”，比如两个积分相等、某个函数可以写成另一个函数的复合。回忆第~\ref{lec:set-systems} 讲的讨论：“好集”的全体总是对差集和递增并封闭，因此是 $\lambda$-类，但一般对交集不封闭。函数的情形完全类似：“好函数”的全体通常对线性组合和递增极限封闭，而对乘积不封闭。下面的定理说，这样的函数类只要包含一个 $\pi$-类上的示性函数，就包含了全体有界可测函数。证明的第一步正是把“好函数”翻译回“好集”，得到一个 $\lambda$-类；然后用简单函数逼近从示性函数走到一般的函数。

\begin{theorem}[函数形式的单调类定理]\label{thm:random-variables:monotone-class}
设 $\cP$ 是 $\Omega$ 上的 $\pi$-类，$\Omega\in\cP$。设 $\cH$ 是一族函数 $\Omega\to\R$，满足
\begin{enumerate}
\item 对一切 $A\in\cP$，$\one_A\in\cH$；
\item 若 $f,g\in\cH$，$c\in\R$，则 $f+g\in\cH$，$cf\in\cH$；
\item 若 $f_n\in\cH$，$f_n\geqslant0$，$f_n\uparrow f$，且 $f$ 有界，则 $f\in\cH$。
\end{enumerate}
则 $\cH$ 包含全体有界 $\sigma(\cP)$-可测函数。
\end{theorem}
\begin{proof}
\emph{第一步：从函数回到集合。}令 $\cL=\{A\subseteq\Omega:\one_A\in\cH\}$。验证 $\cL$ 是 $\lambda$-类：
\begin{itemize}
\item $\Omega\in\cP$，由 (1)，$\one_\Omega\in\cH$，故 $\Omega\in\cL$。
\item 若 $A,B\in\cL$，$A\subseteq B$，则 $\one_{B\setminus A}=\one_B-\one_A=\one_B+(-1)\one_A\in\cH$，由 (2)。
\item 若 $A_n\in\cL$，$A_n\uparrow A$，则 $\one_{A_n}\geqslant0$，$\one_{A_n}\uparrow\one_A$（命题~\ref{prop:set-systems:limsup}(4)(5)），$\one_A$ 有界，由 (3)，$\one_A\in\cH$。
\end{itemize}
由 (1)，$\cP\subseteq\cL$。由 $\pi$-$\lambda$ 定理，$\sigma(\cP)\subseteq\cL$，即对一切 $A\in\sigma(\cP)$ 有 $\one_A\in\cH$。

\emph{第二步：简单函数。}由 (2) 及第一步，一切 $\sigma(\cP)$-可测的简单函数 $\sum_ic_i\one_{A_i}$ 属于 $\cH$。

\emph{第三步：有界非负可测函数。}设 $f$ 是 $\sigma(\cP)$-可测的，$0\leqslant f\leqslant M$。由定理~\ref{thm:random-variables:simple-approx}（对 $\sigma$-代数 $\sigma(\cP)$ 使用），存在 $\sigma(\cP)$-可测的非负简单函数 $f_n\uparrow f$。由第二步 $f_n\in\cH$，由 (3)，$f\in\cH$。

\emph{第四步：有界可测函数。}设 $f$ 是 $\sigma(\cP)$-可测的，$\abs f\leqslant M$。则 $f^+,f^-$ 是有界非负的 $\sigma(\cP)$-可测函数（推论~\ref{cor:random-variables:arithmetic}），由第三步属于 $\cH$，由 (2)，$f=f^++(-1)f^-\in\cH$。
\end{proof}

\begin{remark}
条件 $\Omega\in\cP$ 可以用“常数函数 $1\in\cH$”代替，证明不变。条件 (3) 只要求对\emph{有界}的递增极限封闭，这在应用中很重要：例如 $\cH$ 是“两个有限测度的积分相等”的函数类时，有界性保证积分有限，才能用单调收敛定理。
\end{remark}

作为第一个应用，我们来说明上面“$\sigma(X)$ 是 $X$ 所包含的信息”这种说法的确切含义：一个随机变量 $Y$ 由 $X$ 的信息决定（即 $\sigma(X)$-可测），当且仅当 $Y$ 是 $X$ 的函数。

\begin{theorem}[Doob--Dynkin 引理]\label{thm:random-variables:doob-dynkin}
设 $X:(\Omega,\cF)\to(S,\cS)$ 可测，$Y:\Omega\to\R$。则 $Y$ 是 $\sigma(X)$-可测的，当且仅当存在 $\cS/\cR$-可测函数 $f:S\to\R$ 使 $Y=f\circ X$。
\end{theorem}
\begin{proof}
充分性：由定理~\ref{thm:random-variables:composition}(1)，$f\circ X$ 是 $\sigma(X)/\cR$-可测的（$X$ 是 $\sigma(X)/\cS$-可测的）。

必要性。令 $\cH=\{Y:\Omega\to\R:\ \text{存在}\ \cS\text{-可测的}\ f:S\to\R\ \text{使}\ Y=f\circ X\}$，$\cP=\sigma(X)$（它是 $\sigma$-代数，当然是含 $\Omega$ 的 $\pi$-类）。逐条验证定理~\ref{thm:random-variables:monotone-class} 的条件：
\begin{enumerate}
\item 设 $A\in\sigma(X)$，$A=X^{-1}(B)$，$B\in\cS$。则 $\one_A(\omega)=\one_B(X(\omega))$，而 $\one_B$ 是 $\cS$-可测的。
\item 若 $Y_1=f_1\circ X$，$Y_2=f_2\circ X$，则 $Y_1+Y_2=(f_1+f_2)\circ X$，$cY_1=(cf_1)\circ X$，$f_1+f_2$ 与 $cf_1$ 可测。
\item 设 $Y_n=f_n\circ X\in\cH$，$Y_n\geqslant0$，$Y_n\uparrow Y$，$Y$ 有界。令 $g=\limsup_nf_n$，它是 $S$ 上的广义实值可测函数（定理~\ref{thm:random-variables:limits} 对 $(S,\cS)$ 使用）。令 $E=\{s:\abs{g(s)}<\infty\}\in\cS$，$f=g\one_E$（在 $E^c$ 上令 $f=0$），则 $f$ 是实值 $\cS$-可测函数。对每个 $\omega$，$f_n(X(\omega))=Y_n(\omega)\uparrow Y(\omega)\in\R$，所以 $g(X(\omega))=Y(\omega)$ 有限，$X(\omega)\in E$，$f(X(\omega))=Y(\omega)$。故 $Y\in\cH$。
\end{enumerate}
由定理~\ref{thm:random-variables:monotone-class}，每个有界 $\sigma(X)$-可测函数都属于 $\cH$。

对无界的 $Y$：$Z=\arctan Y$ 是有界 $\sigma(X)$-可测函数（$\arctan$ 连续），所以 $Z=h\circ X$，$h$ 是 $\cS$-可测的。令 $E'=\{s:\abs{h(s)}<\pi/2\}\in\cS$，$f=\tan(h)\one_{E'}$（在 $E'^c$ 上为 $0$）。由于 $\tan$ 在 $(-\pi/2,\pi/2)$ 上连续，$f$ 是 $\cS$-可测的（在 $E'$ 上是连续函数与可测函数的复合，见习题~\exref{ex:random-variables:A1}）。对每个 $\omega$，$h(X(\omega))=\arctan Y(\omega)\in(-\pi/2,\pi/2)$，所以 $f(X(\omega))=\tan(\arctan Y(\omega))=Y(\omega)$。
\end{proof}

\begin{remark}
Doob--Dynkin 引理的两个方向难易程度很不一样。“$Y=f(X)$ 蕴含 $Y$ 是 $\sigma(X)$-可测的”只是复合函数的可测性；反过来，由 $\sigma(X)$-可测性造出 $f$ 才是实质性的内容。直观上，若 $X(\omega)=X(\omega')$，则 $\omega$ 与 $\omega'$ 属于 $\sigma(X)$ 中完全相同的集合，一个 $\sigma(X)$-可测的 $Y$ 无法区分它们，只能取相同的值，于是可以令 $f(X(\omega))=Y(\omega)$。证明的难处在于说明这样定义的 $f$ 是可测的，并且在 $X(\Omega)$ 之外也要合理地定义，单调类定理帮我们绕开了这些细节。
\end{remark}

\begin{example}
设 $X=\one_A$。则 $\sigma(X)$-可测的随机变量恰好是形如 $a\one_A+b\one_{A^c}$ 的函数，即 $f(X)$，其中 $f(1)=a$，$f(0)=b$。设 $(X_1,\dots,X_n)$ 是随机向量，则 $\sigma(X_1,\dots,X_n)$-可测的随机变量恰好是 $f(X_1,\dots,X_n)$，$f$ 为 Borel 函数。以后条件期望 $E(Y\mid X_1,\dots,X_n)$ 就是这样一个函数。
\end{example}

\topic{本讲小结}

\begin{enumerate}
\item 随机变量是可测函数 $X:\Omega\to\R$：每个 $\{X\in B\}$（$B\in\cR$）都是事件。定义中没有概率测度。只需在生成元上验证，比如对一切 $x$ 验证 $\{X\leqslant x\}\in\cF$（推论~\ref{cor:random-variables:criterion}）。
\item 可测性对一切常见运算封闭：复合、连续函数、四则运算、可数个的上下确界与上下极限（定理~\ref{thm:random-variables:composition}、推论~\ref{cor:random-variables:arithmetic}、定理~\ref{thm:random-variables:limits}）。$\{X_n\ \text{收敛}\}$ 是事件，几乎必然收敛有意义。
\item 不可数个随机变量的上确界可以不可测；路径右连续时可以化为可数个（命题~\ref{prop:random-variables:sup-right-continuous}），经验分布函数的 $D_n$ 是随机变量。
\item 乘积 $\sigma$-代数由可测矩形生成，$\cR^d=\cR\otimes\cdots\otimes\cR$；可测集与可测函数的截口可测（定理~\ref{thm:random-variables:sections}），这是第~\ref{lec:product-measures} 讲 Fubini 定理的第一步。
\item $\sigma(X)$ 是 $X$ 所含的信息；$Y$ 是 $\sigma(X)$-可测的当且仅当 $Y=f(X)$（Doob--Dynkin 引理）。
\item 非负可测函数是递增简单函数列的极限（定理~\ref{thm:random-variables:simple-approx}），这是第~\ref{lec:integration} 讲定义积分的出发点；函数形式的单调类定理把 $\pi$-$\lambda$ 定理推广到函数上。
\end{enumerate}

% =====================================================================
\homework{Doob--Dynkin 引理的另一证明，间断点集的可测性，Borel 函数与连续函数的极限}[Doob–Dynkin 引理的另一证明，间断点集，Baire 函数]
\label{hw:random-variables}

\exercisepart{基本习题}

\begin{exercises}{A}{课堂内容的补充}
\item\label{ex:random-variables:A1} 拼接可测函数。
\begin{questions}
\item 设 $X,Y$ 是 $(\Omega,\cF)$ 上的随机变量，$A\in\cF$。证明 $Z=X\one_A+Y\one_{A^c}$（即 $A$ 上等于 $X$，$A^c$ 上等于 $Y$）是随机变量。
\begin{solution}
对任意 $x\in\R$，$Z(\omega)\leqslant x$ 当且仅当“$\omega\in A$ 且 $X(\omega)\leqslant x$”或“$\omega\in A^c$ 且 $Y(\omega)\leqslant x$”，所以
\[
\{Z\leqslant x\}=\bigl(\{X\leqslant x\}\cap A\bigr)\cup\bigl(\{Y\leqslant x\}\cap A^c\bigr)\in\cF .
\]
由推论~\ref{cor:random-variables:criterion}，$Z$ 是随机变量。
\end{solution}
\item 设 $E\in\cS$，$f:S\to\R$，$f|_E$ 对 $E$ 上的迹 $\sigma$-代数可测，证明 $f\one_E$ 是 $\cS$-可测的。
\begin{solution}
$f|_E$ 对迹 $\sigma$-代数 $\cS\cap E=\{C\cap E:C\in\cS\}$ 可测，意思是 $\{s\in E:f(s)\leqslant x\}=C_x\cap E$ 对某个 $C_x\in\cS$ 成立。由于 $E\in\cS$，$C_x\cap E\in\cS$。于是
\[
\{f\one_E\leqslant x\}=\{s\in E:f(s)\leqslant x\}\cup\bigl(E^c\cap\{0\leqslant x\}\bigr)\in\cS,
\]
其中第二部分在 $x\geqslant0$ 时是 $E^c$，否则是空集。
\end{solution}
\end{questions}

\item\label{ex:random-variables:A2} 设 $\cA$ 生成 $\cS$。证明 $X^{-1}(\cA)=\{\{X\in A\}:A\in\cA\}$ 生成 $\sigma(X)$。（这是命题~\ref{prop:random-variables:sigma-generators}，请不看证明独立写一遍。）
\begin{solution}
记 $\cG=\sigma(X^{-1}(\cA))$。每个 $\{X\in A\}$ 属于 $\sigma(X)$，所以 $\cG\subseteq\sigma(X)$。反之，令 $\cB=\{B\subseteq S:\ \{X\in B\}\in\cG\}$，它是 $\sigma$-代数（原像保持补集与可数并），并且包含 $\cA$，所以包含 $\sigma(\cA)=\cS$。这说明对一切 $B\in\cS$，$\{X\in B\}\in\cG$，即 $\sigma(X)\subseteq\cG$。（这也是第~\ref{lec:set-systems} 讲习题~\exref{ex:set-systems:A7} 第 3 小问的特例。）
\end{solution}

\item\label{ex:random-variables:A3} 设 $f:\R\to\R$ 连续，$X_n\to X$ a.s.。证明 $f(X_n)\to f(X)$ a.s.
\begin{solution}
令 $\Omega_0=\{X_n\to X\}$，$P(\Omega_0)=1$。对 $\omega\in\Omega_0$，数列 $X_n(\omega)\to X(\omega)$，由 $f$ 的连续性 $f(X_n(\omega))\to f(X(\omega))$。所以 $\{f(X_n)\to f(X)\}\supseteq\Omega_0$。左边是事件（$f(X_n)$、$f(X)$ 是随机变量，定义~\ref{def:random-variables:as} 后的说明），概率不小于 $P(\Omega_0)=1$。
\end{solution}

\item\label{ex:random-variables:A4} 称 $f:\R^d\to\R$ \term{下半连续}（lower semicontinuous），若对一切 $x$，$\liminf_{y\to x}f(y)\geqslant f(x)$。证明 $f$ 下半连续当且仅当对一切 $a$，$\{x:f(x)\leqslant a\}$ 是闭集，从而下半连续函数是 Borel 函数。
\begin{solution}
按定义，$\liminf_{y\to x}f(y)\geqslant f(x)$ 的意思是：对任意 $\varepsilon>0$，存在 $\delta>0$，使 $\abs{y-x}<\delta$ 时 $f(y)>f(x)-\varepsilon$。

“$\Rightarrow$”：设 $x_k\in\{f\leqslant a\}$，$x_k\to x$。对任意 $\varepsilon>0$，取 $\delta$ 如上，$k$ 充分大时 $\abs{x_k-x}<\delta$，于是 $f(x)-\varepsilon<f(x_k)\leqslant a$。$\varepsilon$ 任意，$f(x)\leqslant a$。所以 $\{f\leqslant a\}$ 是闭集。

“$\Leftarrow$”：若 $f$ 在某点 $x$ 不是下半连续的，则存在 $\varepsilon>0$ 与 $y_k\to x$ 使 $f(y_k)\leqslant f(x)-\varepsilon$。令 $a=f(x)-\varepsilon$，则 $y_k\in\{f\leqslant a\}$，这是闭集，所以 $x\in\{f\leqslant a\}$，即 $f(x)\leqslant f(x)-\varepsilon$，矛盾。

可测性：$\{f\leqslant a\}$ 是闭集，属于 $\cR^d$，由推论~\ref{cor:random-variables:criterion}。
\end{solution}

\item\label{ex:random-variables:A5} 设 $X_n$ 是随机变量。
\begin{questions}
\item 证明 $\{\omega:\sum_nX_n(\omega)\ \text{收敛}\}\in\cF$。
\begin{solution}
部分和 $S_n=X_1+\cdots+X_n$ 是随机变量（推论~\ref{cor:random-variables:arithmetic}）。级数收敛就是 $\lim_nS_n$ 存在且有限，由定理~\ref{thm:random-variables:limits}（对 $S_n$ 用），这个集合属于 $\cF$。也可以直接用 Cauchy 准则写出：
\[
\Bigl\{\sum_nX_n\ \text{收敛}\Bigr\}=\bigcap_{k\geqslant1}\bigcup_{N\geqslant1}\bigcap_{m,n\geqslant N}\Bigl\{\abs{S_n-S_m}\leqslant\frac1k\Bigr\}.
\]
\end{solution}
\item 证明 $\{\omega:\sup_n\abs{X_n(\omega)}<\infty\}\in\cF$。
\begin{solution}
$\sup_n\abs{X_n}$ 是广义实值随机变量（定理~\ref{thm:random-variables:limits}），所以 $\{\sup_n\abs{X_n}<\infty\}=\bigcup_{m\geqslant1}\{\sup_n\abs{X_n}\leqslant m\}\in\cF$。
\end{solution}
\end{questions}

\item\label{ex:random-variables:A6} 在 $([-1,1],\cR\cap[-1,1])$ 上令 $X(\omega)=\omega^2$。证明 $\sigma(X)=\{B\in\cR\cap[-1,1]: B=-B\}$。
\begin{solution}
记右边为 $\cE$。

“$\subseteq$”：$X$ 连续，所以 $\{X\in C\}$ 是 Borel 集；又 $\omega\in\{X\in C\}\iff\omega^2\in C\iff(-\omega)^2\in C$，所以它关于原点对称。

“$\supseteq$”：设 $B\in\cE$。令 $C=\{\omega^2:\omega\in B\cap[0,1]\}$。映射 $h(\omega)=\omega^2$ 是 $[0,1]$ 到自身的连续双射，逆映射 $\sqrt{\cdot}$ 也连续，所以 $C=h(B\cap[0,1])=(\sqrt{\cdot})^{-1}(B\cap[0,1])$ 是 Borel 集。再验证 $\{X\in C\}=B$：$\omega\in\{X\in C\}\iff\omega^2\in C\iff\abs\omega\in B\cap[0,1]\iff\omega\in B$，最后一步用了 $B=-B$。所以 $B\in\sigma(X)$。
\end{solution}

\item\label{ex:random-variables:A7} Borel 函数与 Lebesgue 可测函数。
\begin{questions}
\item 证明 Borel 函数 $f:\R\to\R$ 的全体对逐点极限（极限处处存在且有限时）封闭。
\begin{solution}
设 $f_n$ 是 Borel 函数，$f_n(x)\to f(x)\in\R$ 对每个 $x$ 成立。把定理~\ref{thm:random-variables:limits} 用于可测空间 $(\R,\cR)$ 上的“随机变量” $f_n$：$f=\lim_nf_n=\limsup_nf_n$ 是 $\cR$-可测的。
\end{solution}
\item 举例说明两个 Lebesgue 可测函数（对 Lebesgue 可测集的 $\sigma$-代数可测）的复合不一定是 Lebesgue 可测的。（可以在做完思考题 E 之后再回答。）
\begin{solution}
用问题 E 的记号：$g:[0,1]\to[0,2]$ 是严格递增的连续双射，$V'\subseteq g(C)$ 不是 Lebesgue 可测的，$M=g^{-1}(V')\subseteq C$ 是 Lebesgue 可测集（零测集的子集）。令 $h:\R\to[0,1]$ 在 $[0,2]$ 上等于 $g^{-1}$，在 $(-\infty,0)$ 上等于 $0$，在 $(2,\infty)$ 上等于 $1$，它连续，从而 Lebesgue 可测。令 $\varphi=\one_M$，它是 Lebesgue 可测函数。但
\[
\{\varphi\circ h=1\}\cap[0,2]=\{y\in[0,2]:\ g^{-1}(y)\in M\}=g(M)=V'
\]
不是 Lebesgue 可测的，所以 $\varphi\circ h$ 不是 Lebesgue 可测函数。

问题出在：$\varphi$ 的可测性是对 Lebesgue $\sigma$-代数说的，而 $h$ 只保证 \emph{Borel} 集的原像是 Lebesgue 可测的；Lebesgue 可测集 $M$ 的原像就不一定了。这又一次说明 Borel $\sigma$-代数在复合运算下更好用。
\end{solution}
\end{questions}

\item\label{ex:random-variables:A8} 设 $X,Y$ 是随机变量。证明 $\{X=Y\}$，$\{X<Y\}$ 属于 $\cF$。
\begin{solution}
$X-Y$ 是随机变量（推论~\ref{cor:random-variables:arithmetic}），$\{X=Y\}=\{X-Y\in\{0\}\}$，$\{X<Y\}=\{X-Y\in(-\infty,0)\}$，$\{0\}$ 与 $(-\infty,0)$ 是 Borel 集。也可以直接写：$\{X<Y\}=\bigcup_{q\in\Q}\{X<q\}\cap\{q<Y\}$。
\end{solution}

\item\label{ex:random-variables:A9} 证明 $\R^d$ 上的 Borel $\sigma$-代数 $\cR^d$ 是使一切连续函数 $f:\R^d\to\R$ 可测的最小 $\sigma$-代数。
\begin{solution}
记 $\cG=\sigma(f:\ f\ \text{连续})=\sigma\bigl(\bigcup_f\sigma(f)\bigr)$。连续函数是 Borel 函数（定理~\ref{thm:random-variables:composition}(2)），所以 $\sigma(f)\subseteq\cR^d$，$\cG\subseteq\cR^d$。

反之，设 $U\subseteq\R^d$ 是开集，$U\neq\R^d$。令 $f(x)=\inf_{y\in U^c}\abs{x-y}$ 为 $x$ 到闭集 $U^c$ 的距离，它是连续函数（第~\ref{lec:measures} 讲作业 C1)），并且 $f(x)>0\iff x\notin U^c\iff x\in U$。所以 $U=\{f>0\}\in\sigma(f)\subseteq\cG$。$U=\R^d$ 当然也在 $\cG$ 中。于是 $\cG$ 包含一切开集，$\cR^d\subseteq\cG$。
\end{solution}

\item\label{ex:random-variables:A10} 设 $(S,\cS)$，$(T,\cT)$ 是可测空间，$\pi_S,\pi_T$ 是 $S\times T$ 到两个分量的投影。
\begin{questions}
\item 证明 $\cS\otimes\cT=\sigma(\pi_S,\pi_T)$，即它是使两个投影都可测的最小 $\sigma$-代数。
\begin{solution}
$\sigma(\pi_S)=\{A\times T:A\in\cS\}$，$\sigma(\pi_T)=\{S\times B:B\in\cT\}$。这些都是可测矩形，所以 $\sigma(\pi_S,\pi_T)\subseteq\cS\otimes\cT$。反之，$A\times B=(A\times T)\cap(S\times B)\in\sigma(\pi_S,\pi_T)$，所以一切可测矩形都在其中，$\cS\otimes\cT\subseteq\sigma(\pi_S,\pi_T)$。
\end{solution}
\item 设 $f:S\to\R$ 是 $\cS$-可测的，$g:T\to\R$ 是 $\cT$-可测的，证明 $(s,t)\mapsto f(s)g(t)$ 与 $(s,t)\mapsto f(s)+g(t)$ 都是 $\cS\otimes\cT$-可测的。
\begin{solution}
由 A10-1)，$\pi_S$ 是 $\cS\otimes\cT/\cS$-可测的，所以 $f\circ\pi_S:(s,t)\mapsto f(s)$ 是 $\cS\otimes\cT$-可测的（定理~\ref{thm:random-variables:composition}(1)）；同理 $(s,t)\mapsto g(t)$ 可测。两个可测函数的积与和可测（推论~\ref{cor:random-variables:arithmetic}，它对任意可测空间成立）。
\end{solution}
\end{questions}

\item\label{ex:random-variables:A11} 设 $f:\R\to\R$ 是 Borel 函数。证明 $f$ 的图像 $\{(x,f(x)):x\in\R\}$ 属于 $\cR^2$，并写出它的各个截口。
\begin{solution}
令 $h(x,y)=y-f(x)$。$(x,y)\mapsto f(x)$ 是 Borel 函数 $f$ 与连续投影的复合，$(x,y)\mapsto y$ 连续，所以 $h$ 是 $\cR^2$-可测的。图像 $G=\{h=0\}\in\cR^2$。

截口：$G_x=\{y:(x,y)\in G\}=\{f(x)\}$ 是单点；$G^y=\{x:(x,y)\in G\}=f^{-1}(\{y\})$，是 Borel 集，与定理~\ref{thm:random-variables:sections} 一致。
\end{solution}

\item\label{ex:random-variables:A12} 设 $\{X_t:t\in[0,1]\}$ 是一族随机变量。
\begin{questions}
\item 设对每个 $\omega$，$t\mapsto X_t(\omega)$ 在 $[0,1]$ 上连续。证明 $\sup_{t\in[0,1]}X_t$ 是随机变量。
\begin{solution}
连续函数右连续。取 $T=[0,1]$，$D=\Q\cap[0,1]$：对 $t<1$，可取 $D$ 中从右边趋于 $t$ 的数列；$t=1$ 时 $1\in D$。由命题~\ref{prop:random-variables:sup-right-continuous}，$\sup_{t\in[0,1]}X_t=\sup_{t\in D}X_t$ 是广义实值随机变量。又连续函数在紧集 $[0,1]$ 上有界，上确界有限，所以它是（实值）随机变量。
\end{solution}
\item 在同样的假设下，证明 $\{\omega:\ \text{存在}\ t\in[0,1]\ \text{使}\ X_t(\omega)=0\}\in\cF$。
\begin{solution}
$t\mapsto\abs{X_t(\omega)}$ 连续，在 $[0,1]$ 上取到最小值。所以存在 $t$ 使 $X_t(\omega)=0$，当且仅当 $\min_t\abs{X_t(\omega)}=0$，当且仅当 $\inf_{t\in[0,1]}\abs{X_t(\omega)}=0$。对 $\abs{X_t}$ 用命题~\ref{prop:random-variables:sup-right-continuous}，$\inf_{t\in[0,1]}\abs{X_t}=\inf_{t\in D}\abs{X_t}$ 是随机变量，所求的事件是 $\{\inf_{t\in D}\abs{X_t}\leqslant0\}\in\cF$。
\end{solution}
\item 若只假设 $t\mapsto X_t(\omega)$ 左连续，命题~\ref{prop:random-variables:sup-right-continuous} 应当怎样修改？
\begin{solution}
把“从右边逼近”换成“从左边逼近”：设 $D\subseteq T$ 可数，对每个 $t\in T$ 存在 $d_n\in D$，$d_n\leqslant t$，$d_n\to t$（例如 $T=[a,b]$，$D=(\Q\cap[a,b])\cup\{a\}$）。若 $t\mapsto X_t(\omega)$ 在 $T$ 上左连续，则 $\sup_{t\in T}X_t=\sup_{t\in D}X_t$，$\inf$ 同理。证明与原命题完全相同：$X_t(\omega)=\lim_nX_{d_n}(\omega)\leqslant\sup_{s\in D}X_s(\omega)$。
\end{solution}
\end{questions}

\item\label{ex:random-variables:A13} 在命题~\ref{prop:random-variables:empirical} 中，设 $F$ 连续。把 $X_1(\omega),\dots,X_n(\omega)$ 从小到大排列为 $X_{(1)}(\omega)\leqslant\cdots\leqslant X_{(n)}(\omega)$。证明
\[
D_n=\max_{1\leqslant k\leqslant n}\max\Bigl\{\frac kn-F(X_{(k)}),\ F(X_{(k)})-\frac{k-1}n\Bigr\},
\]
由此再证明一次 $D_n$ 是随机变量。（需要先说明每个 $X_{(k)}$ 是随机变量：$X_{(k)}=\min_{\abs I=k}\max_{i\in I}X_i$，其中 $I$ 取遍 $\{1,\dots,n\}$ 的 $k$ 元子集。）
\begin{solution}
固定 $\omega$，记 $x_{(k)}=X_{(k)}(\omega)$，约定 $x_{(0)}=-\infty$，$x_{(n+1)}=+\infty$，$F(-\infty)=0$，$F(+\infty)=1$。记右边的最大值为 $M$。

\emph{第一步：$D_n\leqslant M$。}区间 $I_k=[x_{(k)},x_{(k+1)})$（$k=0,\dots,n$）把 $\R$ 分成不交的几段（有些可能是空的）。若 $I_k$ 非空，对 $x\in I_k$，恰好有 $k$ 个观测值不超过 $x$，$F_n(x)=\frac kn$。由 $F$ 单调不减且连续，
\[
\sup_{x\in I_k}\Bigl(\frac kn-F(x)\Bigr)=\frac kn-F(x_{(k)}),\qquad\sup_{x\in I_k}\Bigl(F(x)-\frac kn\Bigr)=F(x_{(k+1)})-\frac kn .
\]
当 $1\leqslant k\leqslant n$ 时，第一式是 $M$ 中的项 $\frac kn-F(x_{(k)})$；当 $k=0$ 时，第一式是 $\sup_{x<x_{(1)}}(-F(x))\leqslant0$。当 $0\leqslant k\leqslant n-1$ 时，第二式是 $M$ 中的项 $F(x_{(j)})-\frac{j-1}n$（$j=k+1$）；当 $k=n$ 时，第二式是 $\sup_x(F(x)-1)\leqslant0$。另一方面 $M\geqslant0$：$M$ 中的两项 $\frac1n-F(x_{(1)})$ 与 $F(x_{(1)})-0$ 之和为 $\frac1n$，其中至少有一个是正的。所以对每个非空的 $I_k$，$\sup_{x\in I_k}\abs{F_n(x)-F(x)}\leqslant M$。对 $k$ 取最大值得 $D_n\leqslant M$。

\emph{第二步：$M\leqslant D_n$。}对每个 $k$，在 $x=x_{(k)}$ 处，$F_n(x_{(k)})\geqslant\frac kn$（至少有 $k$ 个值不超过它），所以 $\frac kn-F(x_{(k)})\leqslant F_n(x_{(k)})-F(x_{(k)})\leqslant D_n$。在 $x\uparrow x_{(k)}$ 时（$x<x_{(k)}$），不超过 $x$ 的观测值至多 $k-1$ 个，$F_n(x)\leqslant\frac{k-1}n$，所以 $F(x)-\frac{k-1}n\leqslant F(x)-F_n(x)\leqslant D_n$；令 $x\uparrow x_{(k)}$，由 $F$ 的连续性，$F(x_{(k)})-\frac{k-1}n\leqslant D_n$。

\emph{可测性。}$X_{(k)}$ 是第 $k$ 小的值，它等于一切 $k$ 元子集 $I$ 上最大值中最小的一个：任何 $k$ 个值的最大值不小于第 $k$ 小的值，而取最小的 $k$ 个值时恰好相等。有限个随机变量的最大、最小值是随机变量（推论~\ref{cor:random-variables:arithmetic}），所以 $X_{(k)}$ 是随机变量；$F$ 连续，$F(X_{(k)})$ 是随机变量；$D_n$ 是有限个随机变量的最大值。

这个公式在统计中很有用：计算 $D_n$ 只需要比较 $2n$ 个数，而不必对一切 $x$ 求上确界。
\end{solution}

\item\label{ex:random-variables:A14} （可测函数类的刻画）设 $(\Omega,\cF)$ 是可测空间。
\begin{questions}
\item 证明：$\cF$-可测的实值函数的全体，是包含全体简单函数、并且对逐点极限封闭的最小函数类。
\begin{solution}
记 $\cC$ 为包含全体简单函数、对逐点极限（极限处处存在且有限）封闭的最小函数类（这样的类的交仍是这样的类，全体函数是一个这样的类，所以最小者存在）。

可测函数的全体包含简单函数，并且对逐点极限封闭（定理~\ref{thm:random-variables:limits}），所以 $\cC$ 包含于它。

反之，设 $X$ 可测。由定理~\ref{thm:random-variables:simple-approx}，存在简单函数 $\varphi_n\to X^+$，$\psi_n\to X^-$，于是简单函数 $\varphi_n-\psi_n\to X$ 逐点成立，$X\in\cC$。
\end{solution}
\item 由 A14-1) 再证明一次 Doob--Dynkin 引理（定理~\ref{thm:random-variables:doob-dynkin}）中的必要性。
\begin{solution}
令 $\cK=\{Y:\ Y=f\circ X,\ f\ \text{是}\ \cS\text{-可测的实值函数}\}$。我们验证 $\cK$ 包含一切 $\sigma(X)$-可测的简单函数，并且对逐点极限封闭，然后对可测空间 $(\Omega,\sigma(X))$ 用 A14-1)。

\emph{简单函数。}$Y=\sum_ic_i\one_{A_i}$，$A_i=X^{-1}(B_i)$，$B_i\in\cS$，则 $Y=f\circ X$，$f=\sum_ic_i\one_{B_i}$。

\emph{逐点极限。}设 $Y_n=f_n\circ X\to Y$ 逐点。令 $E=\{s:\ \limsup_nf_n(s)=\liminf_nf_n(s)\in\R\}$，由定理~\ref{thm:random-variables:limits}（对 $(S,\cS)$ 使用），$E\in\cS$；令 $f=\limsup_nf_n\cdot\one_E$（在 $E^c$ 上为 $0$），它是 $\cS$-可测的实值函数。对每个 $\omega$，$f_n(X(\omega))=Y_n(\omega)$ 收敛到 $Y(\omega)$，所以 $X(\omega)\in E$，$f(X(\omega))=Y(\omega)$。

由 A14-1)，$\cK$ 包含一切 $\sigma(X)$-可测函数。注意关键的一步是在 $E^c$ 上把 $f$ 改成 $0$：在 $X(\Omega)$ 之外，$f_n$ 不一定收敛。
\end{solution}
\end{questions}
\end{exercises}

\begin{probleme}{B}{Doob--Dynkin 引理：用简单函数的证明}
本题不用单调类定理，重新证明定理~\ref{thm:random-variables:doob-dynkin}。设 $X:(\Omega,\cF)\to(S,\cS)$ 可测，$Y$ 是 $\sigma(X)$-可测的实值函数。
\begin{questions}
\item 设 $Y=\one_A$，$A\in\sigma(X)$。找出 $f$ 使 $Y=f\circ X$。
\begin{solution}
$A\in\sigma(X)$ 意味着 $A=X^{-1}(B)$，$B\in\cS$。令 $f=\one_B$，则 $f(X(\omega))=\one_B(X(\omega))=\one_{X^{-1}(B)}(\omega)=\one_A(\omega)$。
\end{solution}
\item 设 $Y=\sum_{i=1}^nc_i\one_{A_i}$ 是 $\sigma(X)$-可测简单函数，$A_i\in\sigma(X)$。找出 $f$。$f$ 是否唯一？在 $X(\Omega)$ 上是否唯一？
\begin{solution}
写 $A_i=X^{-1}(B_i)$，令 $f=\sum_ic_i\one_{B_i}$，由 B1) 与线性性 $f\circ X=Y$。

$f$ 一般不唯一：$B_i$ 只在 $X(\Omega)$ 上被 $A_i$ 决定，在 $S\setminus X(\Omega)$ 上可以任意改动（只要保持可测）。比如 $S=\R$，$X\equiv0$，则 $\sigma(X)=\{\varnothing,\Omega\}$，$Y\equiv c$，任何满足 $f(0)=c$ 的 Borel 函数都可以。

在 $X(\Omega)$ 上 $f$ 是唯一的：若 $f\circ X=g\circ X$，则对 $s=X(\omega)$，$f(s)=Y(\omega)=g(s)$。
\end{solution}
\item 设 $Y\geqslant0$。由定理~\ref{thm:random-variables:simple-approx} 取 $\sigma(X)$-可测简单函数 $Y_n\uparrow Y$，$Y_n=f_n\circ X$。说明不能直接令 $f=\lim f_n$，并说明问题出在 $S\setminus X(\Omega)$ 上。
\begin{solution}
在 $X(\Omega)$ 上，$f_n(X(\omega))=Y_n(\omega)\uparrow Y(\omega)$，极限存在且有限。但由 B2)，$f_n$ 在 $S\setminus X(\Omega)$ 上的值是随意的，比如可以是 $(-1)^n$ 或 $n$，这时 $\lim_nf_n$ 不存在或为 $\infty$，不能作为实值函数 $f$。
\end{solution}
\item 令 $f=\sup_nf_n\cdot\one_{\{\sup_nf_n<\infty\}}$。证明 $f$ 是 $\cS$-可测的实值函数，且 $Y=f\circ X$。
\begin{solution}
$\sup_nf_n$ 是 $S$ 上的广义实值可测函数（定理~\ref{thm:random-variables:limits}），$E=\{\sup_nf_n<\infty\}\in\cS$，在 $E$ 上 $\sup_nf_n$ 是实值的；在 $E^c$ 上令 $f=0$。由习题~\exref{ex:random-variables:A1} 的第 2 小问，$f$ 是 $\cS$-可测的实值函数。对 $s=X(\omega)$，$f_n(s)=Y_n(\omega)$ 递增趋于 $Y(\omega)<\infty$，所以 $\sup_nf_n(s)=Y(\omega)$，$s\in E$，$f(s)=Y(\omega)$。

注意这里不必要求 $f_n$ 在 $S$ 上递增；在 $X(\Omega)$ 上它们自动递增，上确界就是极限。
\end{solution}
\item 对一般的 $Y$ 用 $Y=Y^+-Y^-$ 完成证明。
\begin{solution}
$Y^+=\max\{Y,0\}$ 与 $Y^-$ 是非负的 $\sigma(X)$-可测函数，由 B4)，$Y^+=g\circ X$，$Y^-=h\circ X$，$g,h$ 是 $\cS$-可测的实值函数。令 $f=g-h$，则 $Y=f\circ X$。
\end{solution}
\item 设 $S=\R$，$X$ 是随机变量。证明 $Y$ 是 $\sigma(X)$-可测的当且仅当：对任意 $\omega,\omega'$，$X(\omega)=X(\omega')$ 蕴含 $Y(\omega)=Y(\omega')$ 并且 $Y$ 是 $\cF$-可测的。这个说法对吗？若不对，给出反例并说明缺少了什么。
\begin{solution}
“只当”是对的：若 $Y=f(X)$，则 $X(\omega)=X(\omega')$ 时 $Y(\omega)=Y(\omega')$；$\sigma(X)\subseteq\cF$，$Y$ 当然 $\cF$-可测。

“当”是错的。取 $\Omega=\R$，$\cF=2^\R$，$X(\omega)=\omega$。$X$ 是单射，所以第一个条件对任何 $Y$ 都成立；$\cF=2^\R$，所以任何 $Y$ 都是 $\cF$-可测的。但 $\sigma(X)=\cR$，取一个非 Borel 集 $V$（第~\ref{lec:set-systems} 讲作业 F 说明它存在），$Y=\one_V$ 不是 $\sigma(X)$-可测的。

缺少的是：由第一个条件，可以在 $X(\Omega)$ 上定义 $f(X(\omega))=Y(\omega)$，但这个 $f$ 未必可测。$Y$ 的 $\cF$-可测性不能保证 $f$ 的 $\cS$-可测性，因为 $\cF$ 可能比 $\sigma(X)$ 大得多。
\end{solution}
\end{questions}
\end{probleme}

\begin{probleme}{C}{任意函数的间断点集是 Borel 集}
设 $f:\R^d\to\R$ 是\emph{任意}函数（不假设可测）。本题证明 $f$ 的间断点集是 Borel 集。记 $\abs z$ 为欧氏范数。对 $\delta>0$ 令
\[
f^\delta(x)=\sup\{f(y):\abs{y-x}<\delta\},\qquad f_\delta(x)=\inf\{f(y):\abs{y-x}<\delta\},
\]
取值于 $\overline\R$。
\begin{questions}
\item 证明 $f^\delta$ 是下半连续的：对任意 $a$，$\{x:f^\delta(x)>a\}$ 是开集。同理 $f_\delta$ 上半连续。
\begin{solution}
设 $f^\delta(x)>a$。由上确界的定义，存在 $y$ 使 $\abs{y-x}<\delta$，$f(y)>a$。令 $r=\delta-\abs{y-x}>0$。若 $\abs{x'-x}<r$，则 $\abs{y-x'}\leqslant\abs{y-x}+\abs{x-x'}<\delta$，所以 $f^\delta(x')\geqslant f(y)>a$。于是以 $x$ 为心、$r$ 为半径的球含于 $\{f^\delta>a\}$，后者是开集。$f_\delta$ 同理：$\{f_\delta<a\}$ 是开集。
\end{solution}
\item 证明 $\delta\downarrow0$ 时，$f^\delta(x)$ 单调不增，$f_\delta(x)$ 单调不减，从而极限 $f^0(x)=\lim_{\delta\downarrow0}f^\delta(x)$ 与 $f_0(x)=\lim_{\delta\downarrow0}f_\delta(x)$ 存在，并且 $f^0=\lim_{n\to\infty}f^{1/n}$，$f_0=\lim_nf_{1/n}$。
\begin{solution}
$\delta'<\delta$ 时，半径 $\delta'$ 的球含于半径 $\delta$ 的球，较小集合上的上确界较小，所以 $f^{\delta'}(x)\leqslant f^\delta(x)$；下确界相反。单调函数在 $\delta\downarrow0$ 时的极限在 $\overline\R$ 中存在，并且等于沿任何一个趋于 $0$ 的递减数列（如 $\delta=\frac1n$）的极限。
\end{solution}
\item 由 C1)、C2) 与习题 A4 证明 $f^0$、$f_0$ 是（广义实值）Borel 函数。
\begin{solution}
由 C1)，$\{f^\delta\leqslant a\}$ 是闭集，从而对一切实数 $a$ 属于 $\cR^d$，由命题~\ref{prop:random-variables:extended}，$f^\delta$ 是广义实值 Borel 函数。（这正是习题 A4 的论证，只是现在函数可以取 $\pm\infty$。）由 C2)，$f^0=\inf_nf^{1/n}$，由定理~\ref{thm:random-variables:limits} 是广义实值 Borel 函数。同理 $f_0=\sup_nf_{1/n}$。
\end{solution}
\item 证明 $f$ 在 $x$ 连续，当且仅当 $f^0(x)=f_0(x)$（此时二者有限）。
\begin{solution}
总有 $f_\delta(x)\leqslant f(x)\leqslant f^\delta(x)$（$y=x$ 在球内），所以 $f_0(x)\leqslant f(x)\leqslant f^0(x)$。

“$\Rightarrow$”：对 $\varepsilon>0$，取 $\delta$ 使 $\abs{y-x}<\delta$ 时 $\abs{f(y)-f(x)}<\varepsilon$，则 $f^\delta(x)\leqslant f(x)+\varepsilon$，$f_\delta(x)\geqslant f(x)-\varepsilon$，于是 $0\leqslant f^0(x)-f_0(x)\leqslant2\varepsilon$。$\varepsilon$ 任意，二者相等（都夹在有限数之间，有限）。

“$\Leftarrow$”：若 $f^0(x)=f_0(x)$，则二者都等于 $f(x)$。对 $\varepsilon>0$，由极限的定义取 $\delta$ 使 $f^\delta(x)<f(x)+\varepsilon$，$f_\delta(x)>f(x)-\varepsilon$。于是 $\abs{y-x}<\delta$ 时 $f(x)-\varepsilon<f(y)<f(x)+\varepsilon$，$f$ 在 $x$ 连续。
\end{solution}
\item 由 C3)、C4) 与定理~\ref{thm:random-variables:limits} 的证明方法，证明 $\{x: f\ \text{在}\ x\ \text{处不连续}\}\in\cR^d$。
\begin{solution}
由 C4)，间断点集是 $\{f_0<f^0\}$（总有 $f_0\leqslant f^0$）。与定理~\ref{thm:random-variables:limits} 的证明相同，
\[
\{f_0<f^0\}=\bigcup_{q\in\Q}\bigl(\{f_0<q\}\cap\{f^0>q\}\bigr)\in\cR^d .
\]
这个结论的有趣之处在于 $f$ 本身可以不可测，它的间断点集却总是 Borel 集。由此还可以推出：任何函数的间断点集都是可数个闭集的并（$\{f^0-f_0\geqslant\frac1k\}$ 是闭集）。
\end{solution}
\end{questions}
\end{probleme}

\begin{probleme}{D}{Borel 函数是连续函数的“逐次极限”}
令 $\cB_\ast$ 为满足下列两条的\emph{最小}函数类：(a) 一切连续函数 $\R\to\R$ 属于 $\cB_\ast$；(b) 若 $f_n\in\cB_\ast$ 且 $f_n\to f$ 逐点收敛（$f$ 实值），则 $f\in\cB_\ast$。（$\cB_\ast$ 称为 \term{Baire 函数类}。）本题证明 $\cB_\ast$ 恰好是全体 Borel 函数。
\begin{questions}
\item 说明满足 (a)(b) 的函数类的交仍满足 (a)(b)，从而 $\cB_\ast$ 有定义；并由习题 A7 证明 $\cB_\ast$ 中的函数都是 Borel 函数。
\begin{solution}
若每个类都含连续函数，交也含；若 $f_n$ 属于交，则属于每个类，极限 $f$ 属于每个类，从而属于交。全体函数 $\R\to\R$ 满足 (a)(b)，所以交非空，$\cB_\ast$ 是它们的交。Borel 函数的全体满足 (a)（定理~\ref{thm:random-variables:composition}(2)）与 (b)（习题 A7-1)），所以 $\cB_\ast\subseteq$ Borel 函数。
\end{solution}
\item 固定连续函数 $g$，令 $\cK_g=\{f\in\cB_\ast: f+g\in\cB_\ast\}$。证明 $\cK_g$ 满足 (a)(b)，从而 $\cK_g=\cB_\ast$。
\begin{solution}
(a)：$f$ 连续时 $f+g$ 连续，属于 $\cB_\ast$。(b)：若 $f_n\in\cK_g$，$f_n\to f$，则 $f\in\cB_\ast$，且 $f_n+g\in\cB_\ast$，$f_n+g\to f+g$，所以 $f+g\in\cB_\ast$，$f\in\cK_g$。由 $\cB_\ast$ 的最小性，$\cB_\ast\subseteq\cK_g\subseteq\cB_\ast$。
\end{solution}
\item 固定 $h\in\cB_\ast$，令 $\cK^h=\{f\in\cB_\ast: f+h\in\cB_\ast\}$。由 D2) 证明 $\cK^h$ 包含连续函数，再证明它满足 (b)。从而 $\cB_\ast$ 对加法封闭。同理证明对数乘封闭。
\begin{solution}
若 $f$ 连续，由 D2)（取 $g=f$），$h\in\cK_f$，即 $h+f\in\cB_\ast$，所以 $f\in\cK^h$。(b) 的验证与 D2) 相同。所以 $\cK^h=\cB_\ast$，即对一切 $f,h\in\cB_\ast$，$f+h\in\cB_\ast$。这又是“两次放大”：先让一个加数是连续的，再放开它。

数乘：固定 $c$，$\{f\in\cB_\ast:cf\in\cB_\ast\}$ 含连续函数，对逐点极限封闭，所以是 $\cB_\ast$。
\end{solution}
\item 证明对 $a<b$，$\one_{(a,b]}$ 是一列连续函数的逐点极限，从而属于 $\cB_\ast$。
\begin{solution}
对 $n>\frac1{b-a}$，令 $\varphi_n$ 在 $(-\infty,a]$ 与 $[b+\frac1n,\infty)$ 上为 $0$，在 $[a+\frac1n,b]$ 上为 $1$，在 $[a,a+\frac1n]$ 与 $[b,b+\frac1n]$ 上线性连接。$\varphi_n$ 连续。逐点：$x\leqslant a$ 时 $\varphi_n(x)=0$；$a<x\leqslant b$ 时，$n$ 充分大使 $a+\frac1n\leqslant x$，$\varphi_n(x)=1$；$x>b$ 时，$n$ 充分大使 $b+\frac1n<x$，$\varphi_n(x)=0$。所以 $\varphi_n\to\one_{(a,b]}$。
\end{solution}
\item 由定理~\ref{thm:random-variables:monotone-class}（取 $\cP=\{(a,b]\}\cup\{\varnothing,\R\}$，注意 $\one_\R$ 连续）证明：一切有界 Borel 函数属于 $\cB_\ast$。
\begin{solution}
$\cP$ 是 $\pi$-类，含 $\R$，$\sigma(\cP)=\cR$。取 $\cH=\cB_\ast$，验证定理的三个条件：(1) $\one_\varnothing=0$、$\one_\R=1$ 连续，$\one_{(a,b]}\in\cB_\ast$（D4)）；(2) D3)；(3) 递增有界的极限是逐点极限，由 (b)。所以 $\cB_\ast$ 包含一切有界 $\cR$-可测函数。
\end{solution}
\item 对无界的 Borel 函数 $f$，考虑 $\max\{-n,\min\{f,n\}\}$，完成证明。
\begin{solution}
$f_n=\max\{-n,\min\{f,n\}\}$ 是有界 Borel 函数，由 D5) 属于 $\cB_\ast$；$n\geqslant\abs{f(x)}$ 时 $f_n(x)=f(x)$，所以 $f_n\to f$ 逐点，由 (b)，$f\in\cB_\ast$。结合 D1)，$\cB_\ast$ 恰好是 Borel 函数的全体。

这个结论说明，Borel 函数正是从连续函数出发、反复取逐点极限所能得到的函数。比如 Dirichlet 函数 $\one_\Q=\lim_m\lim_n(\cos(m!\pi x))^{2n}$ 需要取两次极限。
\end{solution}
\end{questions}
\end{probleme}

\thinkingpart

\begin{probleme}*[问题]{E}{一个具体的不是 Borel 集的 Lebesgue 可测集}
设 $C$ 是 Cantor 集（习题~\exref{ex:extension:A4}）。定义 \term{Cantor 函数}（Cantor function）$c:[0,1]\to[0,1]$：若 $x=\sum_n\varepsilon_n3^{-n}\in C$，$\varepsilon_n\in\{0,2\}$，令 $c(x)=\sum_n(\varepsilon_n/2)2^{-n}$；在 $[0,1]\setminus C$ 的每个区间（去掉的开区间）上，令 $c$ 等于该区间端点处的公共值。
\begin{questions}
\item 证明 $c$ 定义合理（两个端点处的值相等），单调不减，连续，$c(0)=0$，$c(1)=1$。
\begin{solution}
\emph{在 $C$ 上定义合理。}由习题 A4-3) 的证明，$C$ 中每个点的 $\{0,2\}$-展开是唯一的。

\emph{端点值相等。}第 $n$ 步去掉的开区间位于 $C_{n-1}$ 的某个区间 $[s,s+3^{-n+1}]$（$s=\sum_{k<n}\varepsilon_k3^{-k}$）的中间，是 $(s+3^{-n},s+2\cdot3^{-n})$。左端点 $s+3^{-n}=s+\sum_{k>n}2\cdot3^{-k}$ 的展开是 $\varepsilon_1\cdots\varepsilon_{n-1}0222\cdots$，右端点的展开是 $\varepsilon_1\cdots\varepsilon_{n-1}2000\cdots$。记 $t=\sum_{k<n}(\varepsilon_k/2)2^{-k}$，则
\[
c(\text{左端点})=t+\sum_{k>n}2^{-k}=t+2^{-n},\qquad c(\text{右端点})=t+2^{-n}.
\]

\emph{单调。}设 $x<y$ 都在 $C$ 中，第一个不同的数字在第 $k$ 位，则 $\varepsilon_k(x)=0$，$\varepsilon_k(y)=2$（否则 $x>y$），于是 $c(y)-c(x)\geqslant2^{-k}-\sum_{j>k}2^{-j}=0$。在去掉的区间上 $c$ 取端点的值，所以整体单调不减。

\emph{连续。}$c$ 的值域包含 $[0,1]$：任给 $y=\sum_kb_k2^{-k}$（$b_k\in\{0,1\}$），$x=\sum_k2b_k3^{-k}\in C$ 满足 $c(x)=y$。单调函数若有跳跃，值域就会缺一段区间；$c$ 的值域是整个 $[0,1]$，所以没有跳跃，连续。

$c(0)=0$（$0$ 的展开全为 $0$），$c(1)=1$（$1=\sum_k2\cdot3^{-k}$）。
\end{solution}
\item 令 $g(x)=x+c(x)$。证明 $g$ 是 $[0,1]$ 到 $[0,2]$ 的严格递增连续双射，其逆也连续。
\begin{solution}
$x$ 严格递增，$c$ 单调不减，和严格递增；两者连续，和连续。$g(0)=0$，$g(1)=2$，由介值定理 $g$ 取遍 $[0,2]$，严格递增给出单射。逆映射 $g^{-1}:[0,2]\to[0,1]$ 严格递增，值域是整个区间 $[0,1]$，与上一问同理没有跳跃，连续。
\end{solution}
\item 证明 $g$ 把 $[0,1]\setminus C$ 中的每个去掉的区间映为同样长度的区间，从而 $\lambda(g([0,1]\setminus C))=1$，$\lambda(g(C))=1$。
\begin{solution}
在去掉的区间 $(u,v)$ 上，$c\equiv c(u)$，所以 $g(x)=x+c(u)$，$g((u,v))=(u+c(u),v+c(u))$，长度 $v-u$。$[0,1]\setminus C$ 是可数个不交的去掉的开区间之并，总长度 $\lambda([0,1]\setminus C)=1-\lambda(C)=1$。$g$ 是单射，像也两两不交，所以 $\lambda(g([0,1]\setminus C))=\sum(v-u)=1$。$g$ 是双射，$g(C)=[0,2]\setminus g([0,1]\setminus C)$，$\lambda(g(C))=2-1=1$。（$g(C)$ 是紧集的连续像，是紧集，当然是 Borel 集。）
\end{solution}
\item 利用第~\ref{lec:extension-of-measures} 讲作业 D 的方法证明：$g(C)$ 含有一个不是 Lebesgue 可测的子集 $V'$。（提示：先证明任何 $\lambda^*(A)>0$ 的集合 $A$ 都含有不可测子集。）
\begin{solution}
\emph{引理。}若 $A\subseteq[0,1)$，$\lambda^*(A)>0$，则 $A$ 有不是 Lebesgue 可测的子集。设 $V$ 是第~\ref{lec:extension-of-measures} 讲作业 D 中的集合，令 $A_q=A\cap(q\oplus V)$，$q\in\Q\cap[0,1)$，则 $A=\bigcup_qA_q$（作业 D3)）。若每个 $A_q$ 都可测，我们断言 $\lambda(A_q)=0$：$E=A_q$ 可测且 $E\subseteq q\oplus V$，则 $E'=(1-q)\oplus E\subseteq V$ 可测（作业 D5) 的第一步），而 $\{r\oplus E':r\in\Q\cap[0,1)\}$ 两两不交（含于两两不交的 $r\oplus V$），测度都等于 $\lambda(E')$，并都含于 $[0,1)$，所以 $\sum_r\lambda(E')\leqslant1$，$\lambda(E')=0$，$\lambda(E)=0$。于是 $\lambda^*(A)\leqslant\sum_q\lambda(A_q)=0$，矛盾。

\emph{应用。}$\lambda(g(C))=1$，所以 $g(C)\cap[0,1)$ 与 $g(C)\cap[1,2]$ 中有一个测度为正；后者平移 $-1$ 后含于 $[0,1]$，去掉一点不影响。对它用引理（平移回去不改变可测性，命题~\ref{prop:extension:lebesgue}(2) 与作业 C），得到 $g(C)$ 的不可测子集 $V'$。
\end{solution}
\item 令 $M=g^{-1}(V')\subseteq C$。证明 $M$ 是 Lebesgue 可测集（它是零外测度集的子集），但 $M$ 不是 Borel 集。（提示：若 $M$ 是 Borel 集，则 $g(M)=(g^{-1})^{-1}(M)$ 是连续函数下 Borel 集的原像。）
\begin{solution}
$M\subseteq C$，$\lambda^*(M)\leqslant\lambda(C)=0$，由定理~\ref{thm:extension:caratheodory}，$M$ 是 Lebesgue 可测的。若 $M$ 是 Borel 集，由于 $g^{-1}:[0,2]\to[0,1]$ 连续，$\{y:g^{-1}(y)\in M\}=g(M)=V'$ 是 Borel 集，从而 Lebesgue 可测，矛盾。
\end{solution}
\item 由此给出习题 A7 第二部分所要的例子。
\begin{solution}
见习题 A7-2) 的解答：$\varphi=\one_M$ 与 $h$（$g^{-1}$ 的连续延拓）都是 Lebesgue 可测的，而 $\varphi\circ h$ 在 $[0,2]$ 上是 $\one_{V'}$，不可测。
\end{solution}
\end{questions}
\end{probleme}

\exercisepart{课外知识：Egorov 定理}

\begin{probleme}{F}{几乎处处收敛与近一致收敛}
设 $(\Omega,\cF,\mu)$ 是有限测度空间，$X_n,X$ 是实值可测函数，$X_n\to X$ a.e.。本题证明 \term{Egorov 定理}：\emph{对任意 $\varepsilon>0$，存在 $A\in\cF$，$\mu(A^c)<\varepsilon$，使 $X_n\to X$ 在 $A$ 上一致成立。}
\begin{questions}
\item 对 $k,N\geqslant1$ 令 $E_{k,N}=\bigcup_{n\geqslant N}\{\abs{X_n-X}\geqslant1/k\}$。证明 $E_{k,N}\in\cF$，且对固定的 $k$，$N\to\infty$ 时 $E_{k,N}\downarrow E_k$，其中 $E_k\subseteq\{X_n\not\to X\}$。
\begin{solution}
$\abs{X_n-X}$ 可测，所以每个 $\{\abs{X_n-X}\geqslant\frac1k\}\in\cF$，可数并属于 $\cF$。$N$ 增大时并的范围缩小，$E_{k,N}$ 递减；$E_k=\bigcap_NE_{k,N}=\limsup_n\{\abs{X_n-X}\geqslant\frac1k\}$ 是“$\abs{X_n-X}\geqslant\frac1k$ 发生无穷多次”。在 $E_k$ 上 $X_n\not\to X$。
\end{solution}
\item 由测度的从上连续，对每个 $k$ 取 $N_k$ 使 $\mu(E_{k,N_k})<\varepsilon2^{-k}$。令 $A=\bigcap_kE_{k,N_k}^c$。证明 $\mu(A^c)<\varepsilon$。
\begin{solution}
$\mu(E_k)\leqslant\mu(X_n\not\to X)=0$。$\mu$ 有限，由从上连续 $\mu(E_{k,N})\downarrow\mu(E_k)=0$，所以 $N_k$ 存在。$A^c=\bigcup_kE_{k,N_k}$，由次可加性 $\mu(A^c)<\sum_k\varepsilon2^{-k}=\varepsilon$。
\end{solution}
\item 证明 $X_n\to X$ 在 $A$ 上一致成立。
\begin{solution}
设 $\omega\in A$。对每个 $k$，$\omega\notin E_{k,N_k}$，即对一切 $n\geqslant N_k$，$\abs{X_n(\omega)-X(\omega)}<\frac1k$。$N_k$ 与 $\omega$ 无关，所以 $n\geqslant N_k$ 时 $\sup_{\omega\in A}\abs{X_n(\omega)-X(\omega)}\leqslant\frac1k$。这就是一致收敛。
\end{solution}
\item 举例说明 $\mu$ 有限的条件不能去掉（考虑 $(\R,\cR,\lambda)$ 上的 $\one_{[n,n+1]}$），并说明结论不能加强为 $\mu(A^c)=0$（考虑 $([0,1],\lambda)$ 上的 $x^n$）。
\begin{solution}
\emph{无限测度。}$X_n=\one_{[n,n+1]}\to0$ 处处成立。设 $A$ 使 $X_n\to0$ 在 $A$ 上一致，则存在 $N$，$n\geqslant N$ 时 $\sup_A\one_{[n,n+1]}<\frac12$，即 $A\cap[n,n+1]=\varnothing$，于是 $A^c\supseteq[N,\infty)$，$\lambda(A^c)=\infty$。所以对 $\varepsilon=1$ 结论不成立。

\emph{不能要求 $\mu(A^c)=0$。}$x^n\to0$ 对 $x\in[0,1)$ 成立。若 $\lambda(A^c)=0$，则 $A$ 在 $1$ 的任何左邻域 $(1-\delta,1)$ 中都有点（否则 $A^c$ 含一个正测度的区间），于是对每个 $n$，$\sup_{x\in A}x^n=1$，不趋于 $0$。

Egorov 定理说的是：在有限测度空间上，几乎处处收敛“差不多”就是一致收敛。这在以后证明依概率收敛与几乎必然收敛的关系时会用到类似的想法。
\end{solution}
\end{questions}
\end{probleme}
